\section{2023 Mathematical Physics}

\textit{Solve every problem.}

\subsection{Classical Mechanics}

\begin{exercise}
\label{physics:2023:1}
Consider the Lagrangian
\begin{equation}
    L (x, y, \dot {x}, \dot {y}) = \frac {1}{2} \frac {\dot {x} ^ {2} + \dot {y} ^ {2} + 2 (x \dot {y} - y \dot {x})}{x ^ {2} + y ^ {2}} \label{eq:2023:1:1}
\end{equation}

\begin{enumerate}[label=(\alph*)]
    \item Compute the Hamiltonian $H ( x , y , p _ { x } , p _ { y } )$ , and show the final form can be written as
    \begin{equation}
        \frac {1}{2} f (x, y) \left[ \left(p _ {x} - A _ {x} (x, y)\right) ^ {2} + \left(p _ {y} - A _ {y} (x, y)\right) ^ {2} \right] \label{eq:2023:1:2}
    \end{equation}
    for some $f , A _ { x } , A _ { y }$. Find the vector potential $\vec { A }$ and then compute the corresponding magnetic field away from the origin. (Hint: recall that $\vec { B } = \mathrm { curl } \vec { A } = \nabla \times \vec { A }$\footnote{In vector calculus, the curl of a 2D vector field $\mathbf{v} = (v_x, v_y)$ is formally defined as the scalar field $(\nabla \times \mathbf{v})_z = \partial_x v_y - \partial_y v_x$, representing the circulation density at a point. More generally, in differential geometry, the curl is the exterior derivative of a 1-form $v = v_x dx + v_y dy$, resulting in a 2-form $dv = (\partial_x v_y - \partial_y v_x) dx \wedge dy$. In 3D, this 2-form is mapped back to a vector field via the Hodge star operator, which is the standard cross-product form used in Maxwell's equations.}).
    \item Prove that the Lagrangian $L ( x , y , \dot { x } , \dot { y } )$ is invariant under two symmetries: rotations and scale transformations.
    \item Derive the conserved quantities for both symmetries.
    \item Rewrite the Lagrangian in polar coordinates, write down the Euler-Lagrange equations and solve them.
\end{enumerate}
\end{exercise}

\begin{solution}
\textbf{Prerequisites / Core Concepts:}
\begin{itemize}
    \item \textbf{Legendre Transformation:} The process of switching from $(q, \dot{q})$ to $(q, p)$, defined by $H = \sum p_i \dot{q}_i - L$.
    \item \textbf{Canonical Momentum:} $p_i = \frac{\partial L}{\partial \dot{q}_i}$.
    \item \textbf{Noether's Theorem:} A principle stating that every continuous symmetry of the action corresponds to a conserved quantity.
    \item \textbf{Conserved Quantities:} Physical observables $Q$ that remain invariant ($dQ/dt = 0$) during the dynamical evolution of a system. Mathematically, for a system with coordinates $q_a$ and canonical momenta $p_a$, a continuous symmetry transformation $\delta q_a$ generated by a parameter $\epsilon$ corresponds to a conserved Noether charge $Q = \sum_a p_a \frac{\delta q_a}{\delta \epsilon}$\sidenote{For a free particle $L = \frac{1}{2}m\dot{x}^2$, the transformation $x \to x + \epsilon$ yields $\delta x/\delta \epsilon = 1$. The conserved quantity is $Q = p \cdot 1 = m\dot{x}$. Since $\partial L/\partial x = 0$, the Euler-Lagrange equations $\frac{d}{dt}p = \frac{\partial L}{\partial x}$ directly imply $\dot{p} = 0$.}. Typical examples include energy (from time translation symmetry), linear momentum (from spatial translation), and angular momentum (from rotational symmetry).
    \item \textbf{Vector Potential:} A vector field $\vec{A}$ whose curl gives the magnetic field $\vec{B} = \nabla \times \vec{A}$.
\end{itemize}

\textbf{Intuition / Motivation:}
This system describes a particle moving in a space that feels like a "vortex" centered at the origin. Even though there is no physical magnetic field away from the center, the presence of the vector potential $\vec{A}$ (the "swirl") acts like an unseen force that makes the particle keep a constant angular speed. Meanwhile, because the "stiffness" of the system drops as you move further away exactly as fast as your speed increases, the particle naturally expands outwards in a perfect geometric logarithmic spiral.

\textbf{Logical Step-by-Step Breakdown:}
\textbf{Step 1: Compute the canonical momenta and invert for velocities.}
First, we find the canonical momenta by taking the partial derivatives of the Lagrangian $L$ with respect to the velocities $\dot{x}$ and $\dot{y}$:
\begin{align*}
    p_x &= \frac{\partial L}{\partial \dot{x}} = \frac{\dot{x} - y}{x^2+y^2} \\
    p_y &= \frac{\partial L}{\partial \dot{y}} = \frac{\dot{y} + x}{x^2+y^2}
\end{align*}
Next, we invert these relations to express the velocities $\dot{x}$ and $\dot{y}$ in terms of the momenta $p_x$ and $p_y$\sidenote{From $p_x(x^2+y^2) = \dot{x} - y$, we isolate $\dot{x} = (x^2+y^2)p_x + y$. Similarly, from $p_y(x^2+y^2) = \dot{y} + x$, we obtain $\dot{y} = (x^2+y^2)p_y - x$. These expressions allow us to eliminate velocities in favor of phase-space variables.}:
\begin{align*}
    \dot{x} &= (x^2+y^2)p_x + y \\
    \dot{y} &= (x^2+y^2)p_y - x
\end{align*}

\textbf{Step 2: Construct the Hamiltonian.}
Using the Legendre transformation $H = p_x \dot{x} + p_y \dot{y} - L$, we substitute the expressions for $\dot{x}$ and $\dot{y}$.
First, note that the Lagrangian is a homogeneous quadratic plus a linear term in velocities:
\begin{align*}
    L &= \frac{\dot{x}^2 + \dot{y}^2}{2(x^2+y^2)} + \frac{x\dot{y} - y\dot{x}}{x^2+y^2}
\end{align*}
By Euler's theorem on homogeneous functions, $p_x \dot{x} + p_y \dot{y} = 2 \left( \frac{\dot{x}^2 + \dot{y}^2}{2(x^2+y^2)} \right) + \frac{x\dot{y} - y\dot{x}}{x^2+y^2}$\sidenote{For a function $f(\dot{q})$ of degree $n$, $\sum \dot{q}_i \frac{\partial f}{\partial \dot{q}_i} = n f$. Here, $L = T_2 + T_1$, where $T_2 = \frac{\dot{x}^2+\dot{y}^2}{2r^2}$ is degree 2 and $T_1 = \frac{x\dot{y}-y\dot{x}}{r^2}$ is degree 1. Thus $\sum \dot{q}_i p_i = 2 T_2 + 1 T_1$.}.
Thus, subtracting $L$:
\begin{equation*}
    H = p_x \dot{x} + p_y \dot{y} - L = \frac{\dot{x}^2 + \dot{y}^2}{2(x^2+y^2)}
\end{equation*}
Now we express this entirely in terms of momenta:
\begin{align*}
    H &= \frac{1}{2(x^2+y^2)} \left[ \left( (x^2+y^2)p_x + y \right)^2 + \left( (x^2+y^2)p_y - x \right)^2 \right] \\
      &= \frac{1}{2}(x^2+y^2) \left[ \left( p_x + \frac{y}{x^2+y^2} \right)^2 + \left( p_y - \frac{x}{x^2+y^2} \right)^2 \right]
\end{align*}
Comparing this to the standard form $\frac{1}{2} f(x, y) \left[ (p_x - A_x)^2 + (p_y - A_y)^2 \right]$, we identify $f(x, y) = x^2+y^2$ and the vector potential $\vec{A} = (A_x, A_y) = \left( \frac{-y}{x^2+y^2}, \frac{x}{x^2+y^2} \right)$.

\textbf{Step 3: Calculate the magnetic field.}
The magnetic field is the curl of the vector potential $\vec{A}$. In 2D, this is the $z$-component:
\begin{align*}
    B_z &= \partial_x A_y - \partial_y A_x \\
        &= \frac{\partial}{\partial x} \left( \frac{x}{x^2+y^2} \right) - \frac{\partial}{\partial y} \left( \frac{-y}{x^2+y^2} \right) \\
        &= \frac{1 \cdot (x^2+y^2) - x \cdot 2x}{(x^2+y^2)^2} - \frac{-1 \cdot (x^2+y^2) - (-y) \cdot 2y}{(x^2+y^2)^2} \\
        &= \frac{y^2-x^2}{(x^2+y^2)^2} - \frac{y^2-x^2}{(x^2+y^2)^2} = 0
\end{align*}
The magnetic field is zero everywhere except at the origin (where it is singular).

\textbf{Step 4: Prove rotational and scale invariance.}
For rotational invariance, we note that $x^2+y^2 = r^2$ and $\dot{x}^2+\dot{y}^2 = v^2$ are both strictly invariant under rotations. The term $x\dot{y} - y\dot{x}$ represents the $z$-component of angular momentum\sidenote{In the $xy$-plane, the angular momentum vector is $\vec{L} = \vec{r} \times \vec{p}$. For a particle of unit mass, $\vec{L} = (x\hat{i} + y\hat{j}) \times (\dot{x}\hat{i} + \dot{y}\hat{j}) = (x\dot{y} - y\dot{x})\hat{k}$. Thus, the term is the projection of the angular momentum onto the axis of rotation. Geometrically, it is invariant under rotation because the cross product magnitude $L = r p \sin\theta$ depends only on the relative orientation and the lengths of the vectors, both of which are rotationally invariant scalars.}, which is also rotationally invariant. Thus, all parts of $L$ are invariant under rotations.
For scale invariance, apply the spatial transformation $x \to \lambda x$ and $y \to \lambda y$. This inherently implies $\dot{x} \to \lambda \dot{x}$ and $\dot{y} \to \lambda \dot{y}$.
The numerator of the Lagrangian scales as:
\begin{equation*}
    (\lambda \dot{x})^2 + (\lambda \dot{y})^2 + 2(\lambda x \lambda \dot{y} - \lambda y \lambda \dot{x}) = \lambda^2 \left( \dot{x}^2 + \dot{y}^2 + 2(x\dot{y}-y\dot{x}) \right)
\end{equation*}
The denominator scales as:
\begin{equation*}
    (\lambda x)^2 + (\lambda y)^2 = \lambda^2 (x^2+y^2)
\end{equation*}
The $\lambda^2$ factors cancel out completely, making the overall Lagrangian $L$ invariant under scale transformations.

\textbf{Step 5: Derive the conserved quantities.}
Using Noether's theorem, we find the conserved quantity associated with each symmetry.
For rotations (parameterized by angle $\theta$), the infinitesimal transformations are $\delta x = -y \delta\theta$ and $\delta y = x \delta\theta$\sidenote{Under a rotation by $\delta\theta$, $x' = x\cos\delta\theta - y\sin\delta\theta \approx x - y\delta\theta$ and $y' = x\sin\delta\theta + y\cos\delta\theta \approx x\delta\theta + y$. Thus $\delta x = -y\delta\theta$ and $\delta y = x\delta\theta$.}. The conserved quantity is the angular momentum:
\begin{equation*}
    J = p_x \frac{\delta x}{\delta \theta} + p_y \frac{\delta y}{\delta \theta} = p_x (-y) + p_y (x) = x p_y - y p_x
\end{equation*}
For scale invariance (parameterized by $\lambda = e^s \approx 1+s$), the infinitesimal transformations are $\delta x = x \delta s$ and $\delta y = y \delta s$\sidenote{A scale transformation $x \to \lambda x$ with $\lambda = 1+\delta s$ gives $x' = (1+\delta s)x = x + x\delta s$. Thus $\delta x = x\delta s$, and similarly $\delta y = y\delta s$.}. The conserved quantity is the dilatation charge:
\begin{equation*}
    D = p_x \frac{\delta x}{\delta s} + p_y \frac{\delta y}{\delta s} = x p_x + y p_y
\end{equation*}

\textbf{Remark: Completeness and Independence of Conserved Quantities}
In a two-dimensional dynamical system, the phase space is four-dimensional $(x, y, p_x, p_y)$. A system is considered "integrable" if it possesses enough independent conserved quantities (integrals of motion) to restrict the motion to a lower-dimensional manifold. For this specific Lagrangian, we have identified two primary symmetries:
\begin{itemize}
    \item \textbf{Rotational Symmetry:} Leads to the angular momentum $J = x p_y - y p_x$.
    \item \textbf{Scale Symmetry:} Leads to the dilatation charge $D = x p_x + y p_y$.
\end{itemize}
One might ask if there are other quantities, such as linear momentum. However, because the Lagrangian explicitly depends on the coordinates $x$ and $y$ through the $x^2+y^2$ denominator, spatial translation symmetry is broken, and linear momentum is \textit{not} conserved. Furthermore, the Hamiltonian $H$ is also conserved because $\partial L / \partial t = 0$, but in this scale-invariant system, $H$ is not functionally independent of $J$ and $D$ in a way that provides new information about the trajectory. These two quantities, $J$ and $D$, are the fundamental "generators" of the symmetries and are sufficient to completely determine the solution in Step 6.

\textbf{Remark: Symmetries vs. Conservation --- How Many Quantities?}
As a student of physics, your intuition is excellent: Part (b) identifies \textit{two} spatial symmetries, and Noether's theorem indeed guarantees that we will find \textit{at least} two conserved quantities. However, in the study of dynamical systems, we distinguish between the \textit{total number} of constants and the number of \textit{functionally independent} ones.
\begin{itemize}
    \item \textbf{"At Least Two":} You are correct that each independent symmetry generator (rotation and scaling) yields its own Noether charge ($J$ and $D$). Because these two transformations "move" the system in different ways, the resulting quantities are independent.
    \item \textbf{"Only Two?":} For a system with $N$ degrees of freedom (here $N=2$ for $x$ and $y$), we can have at most $2N-1 = 3$ independent constants that define the trajectory in phase space. Besides $J$ and $D$, we also have the Hamiltonian $H$ (energy). However, because of the specific "mathematical cleanliness" of this Lagrangian, $H$ is actually a function of the others\sidenote{Specifically, in terms of the Noether charges, the energy is $H = \frac{1}{2} D^2 + \frac{1}{2} (J-1)^2$. Since $H$ can be perfectly predicted if you know $J$ and $D$, it does not provide "new" information for solving the path.}.
\end{itemize}
Thus, while there are technically multiple conserved quantities, there are only \textbf{two fundamental ones} arising from the spatial symmetries of the problem. These two are sufficient to "reduce" our 2D problem into two simple 1D problems, which is why we stop our search here.

\textbf{Step 6: Rewrite the Lagrangian in polar coordinates and solve the equations of motion.}
Substitute the polar transformation $x = r\cos\phi$ and $y = r\sin\phi$. This gives $x^2+y^2 = r^2$, $\dot{x}^2+\dot{y}^2 = \dot{r}^2 + r^2\dot{\phi}^2$, and $x\dot{y} - y\dot{x} = r^2\dot{\phi}$\sidenote{Differentiating the polar coordinates yields $\dot{x} = \dot{r}\cos\phi - r\dot{\phi}\sin\phi$ and $\dot{y} = \dot{r}\sin\phi + r\dot{\phi}\cos\phi$. Then $x\dot{y} - y\dot{x} = (r\cos\phi)(\dot{r}\sin\phi + r\dot{\phi}\cos\phi) - (r\sin\phi)(\dot{r}\cos\phi - r\dot{\phi}\sin\phi) = r^2\dot{\phi}(\cos^2\phi + \sin^2\phi) = r^2\dot{\phi}$.}.
Substituting these into the Lagrangian:
\begin{equation*}
    L = \frac{1}{2} \frac{\dot{r}^2 + r^2\dot{\phi}^2 + 2r^2\dot{\phi}}{r^2} = \frac{1}{2} \left( \frac{\dot{r}}{r} \right)^2 + \frac{1}{2} \dot{\phi}^2 + \dot{\phi}
\end{equation*}
The Euler-Lagrange equation for the angular coordinate $\phi$ is:
\begin{equation*}
    \frac{d}{dt} \left( \frac{\partial L}{\partial \dot{\phi}} \right) - \frac{\partial L}{\partial \phi} = 0 \implies \frac{d}{dt} (\dot{\phi} + 1) = 0
\end{equation*}
This means $\dot{\phi} + 1$ is a constant. Let $\dot{\phi} = \omega$, where $\omega$ is a constant angular velocity.
To solve for $r$, let us introduce the substitution $u = \ln r$. Then $\dot{u} = \frac{\dot{r}}{r}$. The Lagrangian becomes $L = \frac{1}{2}\dot{u}^2 + \frac{1}{2}\dot{\phi}^2 + \dot{\phi}$.
The Euler-Lagrange equation for $u$ is:
\begin{equation*}
    \frac{d}{dt} \left( \frac{\partial L}{\partial \dot{u}} \right) - \frac{\partial L}{\partial u} = 0 \implies \frac{d}{dt} (\dot{u}) = 0 \implies \ddot{u} = 0
\end{equation*}
Integrating this yields $u(t) = vt + c$, which translates back to $r$ as $r(t) = e^{vt+c} = r_0 e^{vt}$.
The full trajectory in polar coordinates is therefore $r(t) = r_0 e^{vt}$ and $\phi(t) = \omega t + \phi_0$, which geometrically describes a logarithmic spiral.
\end{solution}

\textbf{Remark: The Deep Architecture of Noether's Theorem}

Noether's Theorem, formulated by Emmy Noether in 1915\footnote{The theorem was motivated by a problem in Einstein's then-new General Relativity regarding the local conservation of energy. Noether's work not only solved this specific problem but fundamentally unified the study of physics by showing that conservation laws are not accidental "facts" of nature, but necessary consequences of the mathematical symmetries underlying spacetime and internal fields.}, is the fundamental bridge between the geometry of the universe and the laws of physics. It states that every continuous, differentiable symmetry of the action of a physical system corresponds to a unique conserved quantity (a Noether charge).

\textbf{1. Mathematical Derivation in Lagrangian Mechanics}
Consider a system described by a Lagrangian $L(q_a, \dot{q}_a, t)$. Suppose the system possesses a continuous symmetry parameterized by a small value $\epsilon$. Under the transformation $q_a \to q_a + \epsilon \delta q_a$, the Lagrangian is said to be invariant if it changes at most by a total time derivative:
\begin{equation*}
    L \to L + \epsilon \frac{d\Lambda}{dt}
\end{equation*}
This ensures that the action $S = \int L dt$ remains invariant (up to boundary terms), leaving the equations of motion unchanged. To find the conserved quantity $Q$, we calculate the variation of $L$ directly:
\begin{equation*}
    \delta L = \sum_a \left( \frac{\partial L}{\partial q_a} \epsilon \delta q_a + \frac{\partial L}{\partial \dot{q}_a} \epsilon \frac{d}{dt}(\delta q_a) \right)
\end{equation*}
Using the Euler-Lagrange equations $\frac{\partial L}{\partial q_a} = \frac{d}{dt} \frac{\partial L}{\partial \dot{q}_a}$ to substitute the first term:
\begin{align*}
    \delta L &= \epsilon \sum_a \left( \left[ \frac{d}{dt} \frac{\partial L}{\partial \dot{q}_a} \right] \delta q_a + \frac{\partial L}{\partial \dot{q}_a} \frac{d}{dt}(\delta q_a) \right) \\
    &= \epsilon \frac{d}{dt} \left( \sum_a \frac{\partial L}{\partial \dot{q}_a} \delta q_a \right)
\end{align*}
Equating this to our symmetry condition $\delta L = \epsilon \frac{d\Lambda}{dt}$, we find:
\begin{equation*}
    \frac{d}{dt} \left( \sum_a p_a \delta q_a - \Lambda \right) = 0 \implies Q = \sum_a p_a \delta q_a - \Lambda
\end{equation*}
where $p_a = \partial L / \partial \dot{q}_a$ is the canonical momentum\sidenote{The subtraction of $\Lambda$ is critical. If the Lagrangian is strictly invariant ($\Lambda=0$), the conserved quantity is simply the momentum projected onto the symmetry direction. If $L$ changes by a total derivative, $\Lambda$ acts as a "correction" to maintain the balance.}.

\textbf{2. Classical Examples and Scenarios}

\textbf{Example A: Spatial Translation $\implies$ Linear Momentum}
Consider a particle in a potential $V(x)$. If the potential is uniform ($V(x) = V_0$), the system looks identical if we shift the particle to $x \to x + \epsilon$. Here, $\delta x = 1$ and $\delta L = 0$ (so $\Lambda = 0$).
The conserved quantity is $Q = p \cdot 1 = m\dot{x}$. This confirms that linear momentum is conserved whenever the environment is spatially homogeneous.

\textbf{Example B: Rotational Symmetry $\implies$ Angular Momentum}
In a central potential $V(r)$, the Lagrangian $L = \frac{1}{2}m(\dot{r}^2 + r^2\dot{\theta}^2) - V(r)$ does not depend on the angle $\theta$. The transformation $\theta \to \theta + \epsilon$ yields $\delta \theta = 1$ and $\delta L = 0$.
The conserved quantity is $Q = p_\theta = \frac{\partial L}{\partial \dot{\theta}} = mr^2\dot{\theta}$. This is the precisely defined angular momentum, which remains constant because space is isotropic (it has no "preferred" direction).

\textbf{Example C: Time Translation $\implies$ Energy (Hamiltonian)}
If the Lagrangian $L(q, \dot{q})$ has no explicit time dependence ($\partial L / \partial t = 0$), the system is invariant under $t \to t + \epsilon$. This transformation is slightly more subtle as it involves shifting the time coordinate\sidenote{Under $t \to t+\epsilon$, we have $q(t) \to q(t-\epsilon) \approx q(t) - \epsilon\dot{q}$. Thus $\delta q = -\dot{q}$. The Lagrangian changes as $\delta L = -\epsilon \frac{dL}{dt}$, so $\Lambda = -L$.}.
The conserved quantity is $Q = \sum p_a (-\dot{q}_a) - (-L) = L - \sum p_a \dot{q}_a = -H$.
This proves that the total energy (the Hamiltonian) is conserved if the laws of physics are "steady" and do not change as time passes.

\textbf{Example D: Gauge Symmetry $\implies$ Electric Charge}
In Quantum Field Theory, if a complex scalar field $\psi$ is invariant under a phase rotation $\psi \to e^{i\epsilon} \psi$, this "internal" symmetry leads to a conserved current $j^\mu$. The total integral of the time-component of this current over all space is the total Electric Charge. Thus, the fact that you cannot "create" or "destroy" net charge is a direct result of the fact that the "phase" of a quantum particle is relative and not absolute\footnote{This is a specific instance of a $U(1)$ gauge symmetry. In the Standard Model, more complex symmetries like $SU(2)$ and $SU(3)$ lead to the conservation of "weak isospin" and "color charge", which dictate the behavior of nuclear forces.}.

\subsection{Quantum Mechanics}

\begin{exercise}
\label{physics:2023:2}
A particle of mass $m$ in 2 dimensions is confined by an isotropic harmonic oscillator potential of frequency $\omega$ , while subject to a weak and anisotropic perturbation of strength $\alpha \ll 1$ . The total Hamiltonian of the particle is
\begin{equation}
    H = H _ {0} + V = \frac {p _ {x} ^ {2}}{2 m} + \frac {p _ {y} ^ {2}}{2 m} + \frac {1}{2} m \omega^ {2} \left(x ^ {2} + y ^ {2}\right) + \alpha m \omega^ {2} x y \label{eq:2023:2:3}
\end{equation}

\begin{enumerate}[label=(\alph*)]
    \item When $\alpha = 0$ , what are the energies and degeneracies of the three lowest-lying unperturbed states?
    \item Use perturbation theory to correct the energies of the above three states to the first order in $\alpha$ .
    \item Find the exact spectrum of $H$ . (Hint: you may want to rotate $x$ and $y$ into a new coordinates)
    \item Check that the perturbative results in part b. are recovered from the exact spectrum.
\end{enumerate}
\end{exercise}

\begin{solution}
\textbf{Prerequisites / Core Concepts:}
\begin{itemize}
    \item \textbf{Exact Spectrum:} The exact spectrum of a Hamiltonian $H$ is the complete set of its energy eigenvalues $\{E_n\}$ and their corresponding degeneracies, obtained by solving the Schrödinger equation $H|\psi_n\rangle = E_n |\psi_n\rangle$ rigorously without any approximations (such as perturbation theory). In complex systems, finding the exact spectrum often requires identifying symmetries or performing coordinate transformations (e.g., rotating into normal modes) that allow the Hamiltonian to be separated into simpler, solvable components.
    \item \textbf{Degeneracy ($g$):} In quantum mechanics, the degeneracy of an energy level is defined as the number of linearly independent quantum states (eigenstates) that correspond to the same energy eigenvalue. A system with $g > 1$ for a given energy is said to be "degenerate." In such cases, any linear combination of these $g$ independent states is also an eigenstate with the same energy level.
    \item \textbf{Harmonic Oscillator in $N$ Dimensions:} A system of $N$ independent 1D harmonic oscillators with the same frequency $\omega$ is described by the energy levels $E_n = (n + N/2)\hbar\omega$, where $n = \sum_{i=1}^N n_i$ is the total quantum number.
    \begin{itemize}
        \item \textbf{In 2D ($N=2$):} The energy levels are $E_n = (n+1)\hbar\omega$. For a given $n$, the degeneracy is $g_n = n+1$, as there are $n+1$ combinations of $(n_x, n_y)$ that sum to $n$ (specifically: $(n, 0), (n-1, 1), \dots, (0, n)$).
        \item \textbf{In $N$ Dimensions:} The total number of states with the same energy $E_n$ is given by the number of ways to distribute $n$ indistinguishable "quanta" into $N$ distinguishable "bins" (dimensions). This is a standard "stars and bars" combinatorial problem, resulting in the degeneracy:
        \begin{equation*}
            g_n = \binom{n + N - 1}{N - 1} = \frac{(n + N - 1)!}{n! (N - 1)!}
        \end{equation*}
    \end{itemize}
    \item \textbf{Perturbation Theory (Time-Independent):} This is a fundamental mathematical framework used in quantum mechanics to find approximate solutions for a complex system that cannot be solved exactly. We begin with a "simple" system (the unperturbed system) described by a Hamiltonian $H_0$, whose eigenvalues $E_n^{(0)}$ and eigenstates $|n^{(0)}\rangle$ are exactly known. We then introduce a "weak" perturbation $V$, such that the total Hamiltonian is $H = H_0 + V$ (often written as $H = H_0 + \lambda V$, where $\lambda \ll 1$ is a small parameter).
    \begin{itemize}
        \item \textbf{The Core Idea:} We assume that the physical properties (energies and states) of the complex system can be expanded as a power series (similar to a Taylor expansion) in terms of the perturbation strength:
        \begin{align*}
            E_n &= E_n^{(0)} + E_n^{(1)} + E_n^{(2)} + \dots \\
            |n\rangle &= |n^{(0)}\rangle + |n^{(1)}\rangle + |n^{(2)}\rangle + \dots
        \end{align*}
        \item \textbf{First-Order Energy Correction:} The first-order correction to the energy is simply the expectation value of the perturbation in the unperturbed state: $E_n^{(1)} = \langle n^{(0)} | V | n^{(0)} \rangle$. Physically, this represents the "average" shift in energy caused by the interaction.
        \item \textbf{Degenerate vs. Non-Degenerate:} If an energy level is unique (non-degenerate), we apply the formula above directly. However, if multiple states share the same energy (degenerate, like the excited states of the 2D oscillator), the perturbation can "mix" these states. In this case, we must construct a matrix of the perturbation $V$ within that specific energy subspace and diagonalize it. The resulting eigenvalues are the true first-order shifts, and this process often "lifts" the degeneracy, splitting one single energy level into several distinct ones.
    \end{itemize}
    \item \textbf{Degenerate Perturbation Theory (Matrix Construction):} When an energy level is $g$-fold degenerate, there are $g$ linearly independent states $\{|\psi_1\rangle, |\psi_2\rangle, \dots, |\psi_g\rangle\}$ that all share the same unperturbed energy $E^{(0)}$. To find how the perturbation $V$ affects this level, we must determine the "interaction" between all pairs of these states.
    \begin{itemize}
        \item \textbf{Building the Matrix:} We construct a $g \times g$ matrix, often called the \textbf{perturbation matrix} or \textbf{secular matrix}, where each element $V_{ij}$ is the "sandwich" (matrix element) of the operator $V$ between the $i$-th bra and the $j$-th ket:
        \begin{equation*}
            \mathbf{V} = \begin{pmatrix} 
                \langle \psi_1 | V | \psi_1 \rangle & \langle \psi_1 | V | \psi_2 \rangle & \dots & \langle \psi_1 | V | \psi_g \rangle \\
                \langle \psi_2 | V | \psi_1 \rangle & \langle \psi_2 | V | \psi_2 \rangle & \dots & \langle \psi_2 | V | \psi_g \rangle \\
                \vdots & \vdots & \ddots & \vdots \\
                \langle \psi_g | V | \psi_1 \rangle & \langle \psi_g | V | \psi_2 \rangle & \dots & \langle \psi_g | V | \psi_g \rangle
            \end{pmatrix}
        \end{equation*}
        \item \textbf{Physical Meaning:} The diagonal elements $V_{ii}$ represent the energy shift of each individual state if they didn't "talk" to each other. The off-diagonal elements $V_{ij}$ represent the "mixing strength" or coupling between different states caused by the perturbation.
        \item \textbf{Diagonalization:} We solve the eigenvalue equation $\det(\mathbf{V} - E^{(1)}\mathbf{I}) = 0$. The resulting eigenvalues $\{E^{(1)}_1, E^{(1)}_2, \dots, E^{(1)}_g\}$ are the true first-order energy corrections. This process mathematically "chooses" the correct linear combinations of the original states that behave as stable eigenstates under the influence of $V$.
    \end{itemize}
    \item \textbf{Normal Modes:} Rotating the coordinate system to decouple interacting variables into independent 1D harmonic oscillators.
\end{itemize}

\textbf{Intuition / Motivation:}
Imagine two identical pendulums or bells. When completely unconnected ($\alpha = 0$), they oscillate at exactly the same pitch (degenerate energy levels). If you add a weak spring connecting them (the $\alpha x y$ perturbation term), they can no longer vibrate entirely independently. The single pitch splits into two distinct modes: one where they swing together in sync, and one where they swing in opposite directions. In quantum mechanics, this "mode splitting" lifts the degeneracy, shifting one energy level up and the other down. By rotating our coordinates $45^\circ$, we find the true independent directions of these collective swings.

\textbf{Logical Step-by-Step Breakdown:}
\textbf{Step 1: Identify the unperturbed spectrum and degeneracies.}
The unperturbed Hamiltonian is exactly a 2D isotropic harmonic oscillator:
\begin{equation*}
    H_0 = \frac{p_x^2}{2m} + \frac{p_y^2}{2m} + \frac{1}{2}m\omega^2 (x^2 + y^2)
\end{equation*}
Its energy levels are $E_{n_x, n_y}^{(0)} = \left(n_x + n_y + 1\right) \hbar\omega$, with states $|n_x, n_y\rangle$.
For the lowest three energy levels:
\begin{itemize}
    \item \textbf{Ground State ($n_x+n_y=0$):} The only state is $|0, 0\rangle$. The energy is $E_0^{(0)} = \hbar\omega$. The degeneracy is $g = 1$.
    \item \textbf{First Excited State ($n_x+n_y=1$):} The states are $|1, 0\rangle$ and $|0, 1\rangle$. The energy is $E_1^{(0)} = 2\hbar\omega$. The degeneracy is $g = 2$.
    \item \textbf{Second Excited State ($n_x+n_y=2$):} The states are $|2, 0\rangle$, $|1, 1\rangle$, and $|0, 2\rangle$. The energy is $E_2^{(0)} = 3\hbar\omega$. The degeneracy is $g = 3$.
\end{itemize}

\textbf{Step 2: Calculate first-order perturbative corrections to the lowest three states.}
The perturbation is $V = \alpha m\omega^2 x y$. To evaluate the matrix elements of this potential efficiently in the basis of harmonic oscillator states $|n_x, n_y\rangle$, we transition from coordinate variables $(x, y)$ to the algebraic formalism of \textbf{Ladder Operators} (also known as raising and lowering operators).

In this formalism, the coordinate $x$ is expressed as a linear combination of the annihilation operator $a_x$ and the creation operator $a_x^\dagger$:
\begin{equation*}
    x = \sqrt{\frac{\hbar}{2m\omega}} (a_x + a_x^\dagger), \quad y = \sqrt{\frac{\hbar}{2m\omega}} (a_y + a_y^\dagger)
\end{equation*}
The operator $a^\dagger$ is called the \textbf{creation (or raising) operator} because it "creates" a quantum of energy, shifting the state up by one level ($a^\dagger |n\rangle = \sqrt{n+1} |n+1\rangle$). Conversely, $a$ is the \textbf{annihilation (or lowering) operator} because it "removes" a quantum, shifting the state down ($a |n\rangle = \sqrt{n} |n-1\rangle$).

Using these operators is significantly more efficient than direct integration in coordinate space. Instead of calculating $\int \psi_n^* x y \psi_m dx dy$, we simply observe how the operators act on the indices of our states. The constants in front of the operators ensure that the commutator $[x, p] = i\hbar$ is preserved and that the unperturbed Hamiltonian takes the diagonal form $H_0 = \sum_i \hbar\omega (a_i^\dagger a_i + 1/2)$.

Substituting these into the perturbation $V = \alpha m\omega^2 x y$:
\begin{equation*}
    V = \alpha m\omega^2 \left( \sqrt{\frac{\hbar}{2m\omega}} (a_x + a_x^\dagger) \right) \left( \sqrt{\frac{\hbar}{2m\omega}} (a_y + a_y^\dagger) \right) = \frac{1}{2} \alpha \hbar\omega (a_x + a_x^\dagger)(a_y + a_y^\dagger)
\end{equation*}
This algebraic form allows us to compute the corrections for each level:

\begin{itemize}
    \item \textbf{Ground State ($g=1$):} The first-order correction is $\Delta E_0 = \langle 0, 0 | V | 0, 0 \rangle = 0$. Since $V \propto (a_x + a_x^\dagger)(a_y + a_y^\dagger)$, every term in the expansion either lowers the ground state (giving zero) or raises it to a state like $|1, 1\rangle$, which is orthogonal to the original $\langle 0, 0 |$.
    \item \textbf{First Excited State ($g=2$):} The degenerate subspace is spanned by $\{|1, 0\rangle, |0, 1\rangle\}$. We construct the $2 \times 2$ matrix $\mathbf{V}$ where the rows and columns correspond to these two states:
    \begin{equation*}
        \mathbf{V} = \begin{pmatrix} 
            \langle 1, 0 | V | 1, 0 \rangle & \langle 1, 0 | V | 0, 1 \rangle \\
            \langle 0, 1 | V | 1, 0 \rangle & \langle 0, 1 | V | 0, 1 \rangle
        \end{pmatrix}
    \end{equation*}
    The diagonal elements vanish because the operator $V$ always changes the "particle count" in both $x$ and $y$ dimensions. For example, $V|1, 0\rangle$ produces a linear combination of states like $|0, 1\rangle, |2, 1\rangle, \dots$, all of which have zero overlap with the original $\langle 1, 0 |$.
    The off-diagonal element is calculated as:
    \begin{align*}
        \langle 1, 0 | V | 0, 1 \rangle &= \frac{1}{2} \alpha \hbar\omega \langle 1|_x \langle 0|_y (a_x + a_x^\dagger)(a_y + a_y^\dagger) |0\rangle_x |1\rangle_y \\
        &= \frac{1}{2} \alpha \hbar\omega \underbrace{\langle 1|_x a_x^\dagger |0\rangle_x}_{1} \underbrace{\langle 0|_y a_y |1\rangle_y}_{1} = \frac{1}{2} \alpha \hbar\omega
    \end{align*}
    By symmetry (Hermiticity), $\langle 0, 1 | V | 1, 0 \rangle = \frac{1}{2} \alpha \hbar\omega$. Thus, the matrix is:
    \begin{equation*}
        \mathbf{V} = \frac{1}{2} \alpha \hbar\omega \begin{pmatrix} 0 & 1 \\ 1 & 0 \end{pmatrix}
    \end{equation*}
    The eigenvalues (energy shifts) are $E^{(1)} = \pm \frac{1}{2} \alpha \hbar\omega$. The corrected energies for the first excited states are $E_1 \approx 2\hbar\omega \pm \frac{1}{2}\alpha\hbar\omega$.
    \item \textbf{Second Excited State ($g=3$):} The subspace is spanned by $\{|2, 0\rangle, |1, 1\rangle, |0, 2\rangle\}$. The $3 \times 3$ matrix structure is:
    \begin{equation*}
        \mathbf{V} = \begin{pmatrix} 
            \langle 2, 0 | V | 2, 0 \rangle & \langle 2, 0 | V | 1, 1 \rangle & \langle 2, 0 | V | 0, 2 \rangle \\
            \langle 1, 1 | V | 2, 0 \rangle & \langle 1, 1 | V | 1, 1 \rangle & \langle 1, 1 | V | 0, 2 \rangle \\
            \langle 0, 2 | V | 2, 0 \rangle & \langle 0, 2 | V | 1, 1 \rangle & \langle 0, 2 | V | 0, 2 \rangle
        \end{pmatrix}
    \end{equation*}
    Again, all diagonal elements $\langle n_x, n_y | V | n_x, n_y \rangle$ are zero because $V$ must change both $n_x$ and $n_y$ by $\pm 1$. Additionally, the transition element between the "ends" of the subspace $\langle 2, 0 | V | 0, 2 \rangle = 0$ because $V$ only contains single-step operators (like $a_x a_y^\dagger$); it cannot change an index by two units (from $0$ to $2$) in a single action.
    The non-zero "couplings" are:
    \begin{align*}
        \langle 2, 0 | V | 1, 1 \rangle &= \frac{1}{2}\alpha\hbar\omega \langle 2|a_x^\dagger|1\rangle \langle 0|a_y|1\rangle = \frac{1}{2}\alpha\hbar\omega (\sqrt{2})(1) = \frac{\sqrt{2}}{2}\alpha\hbar\omega \\
        \langle 0, 2 | V | 1, 1 \rangle &= \frac{1}{2}\alpha\hbar\omega \langle 0|a_x|1\rangle \langle 2|a_y^\dagger|1\rangle = \frac{1}{2}\alpha\hbar\omega (1)(\sqrt{2}) = \frac{\sqrt{2}}{2}\alpha\hbar\omega
    \end{align*}
    Placing these into their respective slots ($V_{12}, V_{21}, V_{23}, V_{32}$), the full matrix becomes:
    \begin{equation*}
        \mathbf{V} = \frac{1}{2} \alpha \hbar\omega \begin{pmatrix} 0 & \sqrt{2} & 0 \\ \sqrt{2} & 0 & \sqrt{2} \\ 0 & \sqrt{2} & 0 \end{pmatrix}
    \end{equation*}
    To find the eigenvalues $\lambda$, we solve $\det(\mathbf{V} - \lambda \mathbf{I}) = 0$:
    \begin{equation*}
        -\lambda (\lambda^2 - 2) - \sqrt{2}(-\lambda\sqrt{2}) = -\lambda^3 + 4\lambda = 0
    \end{equation*}
    The roots are $\lambda = 0, \pm 2$. Thus, the corrections are $0$, $\alpha\hbar\omega$, and $-\alpha\hbar\omega$. The corrected energies are $3\hbar\omega$, $3\hbar\omega + \alpha\hbar\omega$, and $3\hbar\omega - \alpha\hbar\omega$.
\end{itemize}

\textbf{Step 3: Find the exact spectrum by rotating coordinates.}
To solve the Hamiltonian exactly, we need to eliminate the "coupling" term $\alpha m\omega^2 xy$ which prevents the variables from being independent. Since the potential is symmetric with respect to $x$ and $y$, we perform a $45^\circ$ rotation into a new coordinate system $(u, v)$ where these "normal modes" are decoupled:
\begin{equation*}
    u = \frac{x + y}{\sqrt{2}}, \quad v = \frac{x - y}{\sqrt{2}}
\end{equation*}
Inverting these relations gives $x = \frac{u+v}{\sqrt{2}}$ and $y = \frac{u-v}{\sqrt{2}}$. Let us transform each part of the Hamiltonian:
\begin{enumerate}
    \item \textbf{Kinetic Energy:} The transformation is orthogonal, so the sum of squares of the momenta is preserved: $p_x^2 + p_y^2 = p_u^2 + p_v^2$.
    \item \textbf{Unperturbed Potential:} $x^2 + y^2 = \left(\frac{u+v}{\sqrt{2}}\right)^2 + \left(\frac{u-v}{\sqrt{2}}\right)^2 = \frac{u^2 + 2uv + v^2}{2} + \frac{u^2 - 2uv + v^2}{2} = u^2 + v^2$.
    \item \textbf{Perturbation Term:} $xy = \left(\frac{u+v}{\sqrt{2}}\right)\left(\frac{u-v}{\sqrt{2}}\right) = \frac{u^2 - v^2}{2}$.
\end{enumerate}
Substituting these into the total Hamiltonian $H = H_0 + V$:
\begin{align*}
    H &= \left[ \frac{p_u^2 + p_v^2}{2m} \right] + \frac{1}{2}m\omega^2(u^2 + v^2) + \alpha m\omega^2 \left( \frac{u^2 - v^2}{2} \right) \\
      &= \left[ \frac{p_u^2}{2m} + \frac{1}{2}m\omega^2(1+\alpha)u^2 \right] + \left[ \frac{p_v^2}{2m} + \frac{1}{2}m\omega^2(1-\alpha)v^2 \right]
\end{align*}
This is now the sum of two \textbf{independent} 1D harmonic oscillators with modified frequencies:
\begin{equation*}
    \omega_u = \omega\sqrt{1+\alpha}, \quad \omega_v = \omega\sqrt{1-\alpha}
\end{equation*}
The exact energy eigenvalues are the sum of the energies of these two oscillators:
\begin{equation*}
    E_{n_u, n_v} = \left( n_u + \frac{1}{2} \right) \hbar\omega\sqrt{1+\alpha} + \left( n_v + \frac{1}{2} \right) \hbar\omega\sqrt{1-\alpha}
\end{equation*}
where $n_u, n_v = 0, 1, 2, \dots$ are the quantum numbers for the new normal modes.

\textbf{Step 4: Verify the perturbative results using the exact spectrum.}
We now check if our approximate results from Step 2 match the exact results in the limit of a "weak" perturbation ($\alpha \ll 1$). We use the Taylor expansion $(1+z)^{1/2} = 1 + \frac{1}{2}z - \frac{1}{8}z^2 + \dots$ for the frequencies:
\begin{align*}
    E_{n_u, n_v} &\approx \left( n_u + \frac{1}{2} \right) \hbar\omega \left( 1 + \frac{1}{2}\alpha - \frac{1}{8}\alpha^2 \right) + \left( n_v + \frac{1}{2} \right) \hbar\omega \left( 1 - \frac{1}{2}\alpha - \frac{1}{8}\alpha^2 \right) \\
    &= (n_u + n_v + 1)\hbar\omega + \frac{1}{2}\alpha\hbar\omega(n_u - n_v) - \frac{1}{8}\alpha^2\hbar\omega(n_u + n_v + 1)
\end{align*}
Let us compare this term-by-term with Step 2 for each level:
\begin{itemize}
    \item \textbf{Ground State ($n_u=0, n_v=0$):} The formula gives $E \approx \hbar\omega - \frac{1}{8}\alpha^2\hbar\omega$. To first order in $\alpha$, the shift is \textbf{zero}, matching $\Delta E_0 = 0$.
    \item \textbf{First Excited State ($n_u+n_v=1$):}
    \begin{itemize}
        \item For $(1, 0)$: $E \approx 2\hbar\omega + \frac{1}{2}\alpha\hbar\omega$.
        \item For $(0, 1)$: $E \approx 2\hbar\omega - \frac{1}{2}\alpha\hbar\omega$.
    \end{itemize}
    This perfectly reproduces the eigenvalues $E_1^{(0)} \pm \frac{1}{2}\alpha\hbar\omega$ from the $2 \times 2$ matrix.
    \item \textbf{Second Excited State ($n_u+n_v=2$):}
    \begin{itemize}
        \item For $(2, 0)$: $E \approx 3\hbar\omega + \frac{1}{2}\alpha\hbar\omega(2) = 3\hbar\omega + \alpha\hbar\omega$.
        \item For $(1, 1)$: $E \approx 3\hbar\omega + \frac{1}{2}\alpha\hbar\omega(0) = 3\hbar\omega$.
        \item For $(0, 2)$: $E \approx 3\hbar\omega + \frac{1}{2}\alpha\hbar\omega(-2) = 3\hbar\omega - \alpha\hbar\omega$.
    \end{itemize}
    These results are exactly the shifts $\{+\alpha\hbar\omega, 0, -\alpha\hbar\omega\}$ we found by diagonalizing the $3 \times 3$ matrix.
\end{itemize}
The complete agreement confirms that degenerate perturbation theory correctly captures the splitting of energy levels in the weak-coupling regime.

\end{solution}

\subsection{Electromagnetism}

\begin{exercise}
\label{physics:2023:3}
One can express the electric fields $\vec { E }$ and magnetic fields $\vec { B }$ in terms of the scalar and vector potentials, $A ^ { \mu } = \left( \phi , { \vec { A } } \right)$ .

\begin{enumerate}[label=(\alph*)]
    \item Write down the expression of $\vec { E }$ and $\vec { B }$ in terms of $( \phi , { \vec { A } } )$ and show that the result is unchanged under gauge transformation
    \begin{equation}
        \phi \rightarrow \phi + \frac{\partial}{\partial t} f, \quad \vec {A} \rightarrow \vec {A} - \nabla f, \label{eq:2023:3:4}
    \end{equation}
    where $f = f(\vec{x}, t)$ is a scalar function.
    \item Show that two of the 4 Maxwell equations of $\vec { E }$ and $\vec { B }$ are satisfied automatically in terms of $A ^ { \mu } = \left( \phi , { \vec { A } } \right)$ .
    \item Derive the equations for the scalar and vector potentials from the remaining Maxwell equations in Lorentz gauge.
    \begin{equation}
        \frac{1}{c} \partial_t \phi + \nabla \cdot \vec{A} = 0, \label{eq:2023:3:5}
    \end{equation}
    \item Recall that the Green's function of the wave equation
    \begin{equation}
        \left(\frac{1}{c^2} \partial_t^2 - \nabla^2\right) G(t, \vec{r}) = \delta(t) \delta^3(\vec{r}) \label{eq:2023:3:6}
    \end{equation}
    is
    \begin{equation}
        G (t - t _ {0}, \vec {r} - \vec {r} _ {0}) = \frac {\theta (t - t _ {0})}{4 \pi | \vec {r} - \vec {r} _ {0} |} \delta \left(t - t _ {0} - \frac {| \vec {r} - \vec {r} _ {0} |}{c}\right). \label{eq:2023:3:7}
    \end{equation}
    Assume a particle of electric charge $e$ moves with trajectory $\vec { R } ( t )$ with ${ \vec { v } } ( t ) = d \vec { R } ( t ) / d t$ . Use the Green's function to derive the potential $\phi ( \vec { r } , t )$ and $\vec { A } ( \vec { r } , t )$ of this particle at $( \vec { r } , t )$. You may assume that $|\vec{R}(t)| \ll |\vec{r}|$, $|\vec{R}(t)| \ll ct$ and $|\vec{v}(t)| \ll c$ and expand your result up to the order $\mathcal{O}(1/|\vec{r}|)$ and $\mathcal{O}(|\vec{v}(t)|/c)$. This is also called non-relativistic and far-field approximations.
\end{enumerate}
\end{exercise}

\begin{solution}
\textbf{Prerequisites / Core Concepts:}
\begin{itemize}
    \item \textbf{Vector Calculus Operators on $\mathbb{R}^3$:} For a scalar field $\phi$ and vector field $\vec{A}$:
    \begin{itemize}
        \item \textbf{Gradient ($\nabla \phi$):} A vector field $\left( \frac{\partial \phi}{\partial x}, \frac{\partial \phi}{\partial y}, \frac{\partial \phi}{\partial z} \right)$ pointing in the direction of steepest ascent.
        \item \textbf{Divergence ($\nabla \cdot \vec{A}$):} A scalar field $\frac{\partial A_x}{\partial x} + \frac{\partial A_y}{\partial y} + \frac{\partial A_z}{\partial z}$ representing the "source density" or net flux per unit volume.
        \item \textbf{Curl ($\nabla \times \vec{A}$):} A vector field $\left( \frac{\partial A_z}{\partial y} - \frac{\partial A_y}{\partial z}, \frac{\partial A_x}{\partial z} - \frac{\partial A_z}{\partial x}, \frac{\partial A_y}{\partial x} - \frac{\partial A_x}{\partial y} \right)$ representing the local circulation density.
        \item \textbf{Laplacian ($\nabla^2 \phi$):} The divergence of the gradient, $\sum \frac{\partial^2 \phi}{\partial x_i^2}$, characterizing the "curvature" or local average of the field.
    \end{itemize}
    \item \textbf{The d'Alembertian Operator ($\Box$):} The relativistic generalization of the Laplacian for 4D spacetime: $\Box = \frac{1}{c^2} \frac{\partial^2}{\partial t^2} - \nabla^2$. It is the differential operator for the wave equation.
    \item \textbf{Maxwell's Equations (PDE Perspective):}
    \begin{itemize}
        \item \textbf{Gauss's Law ($\nabla \cdot \vec{E} = 4\pi\rho$):} Sources of electric flux are charges.
        \item \textbf{Faraday's Law ($\nabla \times \vec{E} + \frac{1}{c}\partial_t \vec{B} = 0$):} Changing magnetic flux induces a "curled" electric field.
        \item \textbf{Gauss's Law for Magnetism ($\nabla \cdot \vec{B} = 0$):} There are no magnetic monopoles (no point sources for $\vec{B}$).
        \item \textbf{Ampere-Maxwell Law ($\nabla \times \vec{B} - \frac{1}{c}\partial_t \vec{E} = \frac{4\pi}{c}\vec{j}$):} Magnetic fields are "curled" by currents and changing electric flux.
    \end{itemize}
    \item \textbf{Electromagnetic 4-Potential ($A^\mu$):} A vector in 4D spacetime $A^\mu = (\phi, \vec{A})$. In the language of differential geometry, this is a connection on a principal $U(1)$ bundle. The physical fields $\vec{E}$ and $\vec{B}$ are the "curvature" or "field strength" components of this connection.
    \item \textbf{Poincaré's Lemma and Potential Existence:} A mathematical theorem stating that on a contractible manifold (like $\mathbb{R}^3$), a closed form ($d\omega = 0$) is always exact ($\omega = d\eta$). This is the technical reason why $\nabla \cdot \vec{B} = 0$ guarantees $\vec{B} = \nabla \times \vec{A}$.
\end{itemize}

\textbf{Remark: Mathematical Foundations --- The Logic of Potentials and Gauge}
To a mathematician, Maxwell's four equations are naturally split into two pairs. The first pair (Gauss's Magnetism and Faraday's Law) are called the \textit{homogeneous} equations. Notice they have zeros on the right side:
\begin{align}
    \nabla \cdot \vec{B} &= 0 \tag{H1} \\
    \nabla \times \vec{E} + \frac{1}{c} \frac{\partial \vec{B}}{\partial t} &= 0 \tag{H2}
\end{align}
These equations are actually "integrability conditions." 
Equation (H1) tells us that $\vec{B}$ is a solenoidal field. By Poincaré's Lemma, there must exist a vector field $\vec{A}$ (the vector potential) such that $\vec{B} = \nabla \times \vec{A}$. 
If we substitute this $\vec{B}$ into (H2), we get:
\begin{equation*}
    \nabla \times \vec{E} + \frac{1}{c} \frac{\partial}{\partial t} (\nabla \times \vec{A}) = 0 \implies \nabla \times \left( \vec{E} + \frac{1}{c} \frac{\partial \vec{A}}{\partial t} \right) = 0
\end{equation*}
Since the term in the parenthesis has zero curl, it must be the gradient of some scalar function. By convention, we call this function $-\phi$ (the scalar potential):
\begin{equation*}
    \vec{E} + \frac{1}{c} \frac{\partial \vec{A}}{\partial t} = -\nabla \phi \implies \vec{E} = -\nabla \phi - \frac{1}{c} \frac{\partial \vec{A}}{\partial t}
\end{equation*}
This is the "Genesis" of the potentials. However, this definition is not unique. If you add the gradient of any scalar function $f$ to $\vec{A}$, the curl ($\vec{B}$) remains the same because $\nabla \times (\nabla f) = 0$\sidenote{Mathematically, the identity $\nabla \times (\nabla f) = 0$ is a consequence of the equality of mixed partial derivatives (Clairaut's Theorem). For example, the $z$-component $(\nabla \times \nabla f)_z = \frac{\partial}{\partial x} (\frac{\partial f}{\partial y}) - \frac{\partial}{\partial y} (\frac{\partial f}{\partial x}) = f_{yx} - f_{xy} = 0$. Geometrically, this means that a gradient field $\nabla f$ is "irrotational": the work done by the field around any closed loop is zero, so there is no net "swirl" at any point.}. To keep $\vec{E}$ the same, you must simultaneously shift $\phi$ by $-\frac{1}{c} \partial_t f$. This "freedom" to shift $\phi$ and $\vec{A}$ without changing the physical fields $\vec{E}$ and $\vec{B}$ is called \textbf{Gauge Invariance}.

In this exercise, we utilize this freedom to pick a specific "constraint" called the \textbf{Lorentz Gauge}: $\nabla \cdot \vec{A} + \frac{1}{c} \partial_t \phi = 0$. This particular choice is mathematically beautiful because it transforms the coupled, messy Maxwell equations into two independent, identical wave equations (the d'Alembertian equations) for $\phi$ and $\vec{A}$, which are much easier to solve using Green's functions.

\textbf{Intuition / Motivation:}
Think of the scalar and vector potentials as the hidden "blueprints" for the electric and magnetic fields. Just as there are many different ways to write a computer program to achieve the exact same user interface, there are infinitely many valid potentials (a "gauge choice") that produce the exact same physical fields. By carefully choosing our gauge (the Lorentz gauge), we can transform Maxwell's complicated coupled equations into simple, independent wave equations. The Green's function then tells us something very intuitive: if a charge moves, the "news" of that movement ripples outward precisely at the speed of light, so the field we feel here right now depends on what the charge was doing some time ago.

\textbf{Logical Step-by-Step Breakdown:}
\textbf{Step 1: Define physical fields in terms of potentials and prove gauge invariance.}
The electric and magnetic fields are defined by the potentials $A^\mu = (\phi, \vec{A})$ as:
\begin{align*}
    \vec{B} &= \nabla \times \vec{A} \\
    \vec{E} &= -\nabla \phi - \frac{1}{c}\frac{\partial \vec{A}}{\partial t}
\end{align*}
Under a gauge transformation with a scalar function $f(t, \vec{r})$, the new potentials are:
\begin{equation*}
    \phi' = \phi + \frac{1}{c}\frac{\partial f}{\partial t}, \quad \vec{A}' = \vec{A} - \nabla f
\end{equation*}
Let us compute the new magnetic field:
\begin{equation*}
    \vec{B}' = \nabla \times \vec{A}' = \nabla \times (\vec{A} - \nabla f) = \nabla \times \vec{A} - \nabla \times (\nabla f) = \vec{B} - 0 = \vec{B}
\end{equation*}
And the new electric field:
\begin{align*}
    \vec{E}' &= -\nabla \phi' - \frac{1}{c}\frac{\partial \vec{A}'}{\partial t} = -\nabla \left( \phi + \frac{1}{c}\frac{\partial f}{\partial t} \right) - \frac{1}{c}\frac{\partial}{\partial t} \left( \vec{A} - \nabla f \right) \\
    &= -\nabla \phi - \frac{1}{c}\nabla\left(\frac{\partial f}{\partial t}\right) - \frac{1}{c}\frac{\partial \vec{A}}{\partial t} + \frac{1}{c}\frac{\partial}{\partial t}(\nabla f) \\
    &= \left( -\nabla \phi - \frac{1}{c}\frac{\partial \vec{A}}{\partial t} \right) - \frac{1}{c}\nabla\left(\frac{\partial f}{\partial t}\right) + \frac{1}{c}\nabla\left(\frac{\partial f}{\partial t}\right) \\
    &= \vec{E}
\end{align*}
Thus, the physical fields remain perfectly unchanged.

\textbf{Step 2: Show that the homogeneous Maxwell equations are automatically satisfied.}
The two homogeneous Maxwell equations are Gauss's Law for magnetism and Faraday's Law of induction.
For Gauss's law for magnetism, we use the vector identity that the divergence of a curl is always zero\sidenote{Mathematically, $\nabla \cdot (\nabla \times \vec{A}) = 0$ is an identity arising from the equality of mixed partial derivatives (Clairaut's Theorem). In Cartesian components, the sum expands as $\frac{\partial}{\partial x}(\frac{\partial A_z}{\partial y} - \frac{\partial A_y}{\partial z}) + \frac{\partial}{\partial y}(\frac{\partial A_x}{\partial z} - \frac{\partial A_z}{\partial x}) + \frac{\partial}{\partial z}(\frac{\partial A_y}{\partial x} - \frac{\partial A_x}{\partial y})$. Grouping terms with the same vector component (e.g., $A_z$) yields $(A_{z,yx} - A_{z,xy})$, which is zero. Geometrically, this means that a "curled" field has no sources or sinks; it flows exclusively in closed loops, and Thus its net flux through any closed surface is always zero.}:
\begin{equation*}
    \nabla \cdot \vec{B} = \nabla \cdot (\nabla \times \vec{A}) = 0
\end{equation*}
For Faraday's law, we take the curl of the electric field. We use the identity that the curl of a gradient is always zero\sidenote{Mathematically, the identity $\nabla \times (\nabla \phi) = 0$ follows from the fact that for a smooth scalar field, the order of differentiation does not matter (Clairaut's Theorem). In the $z$-component of the curl, $(\nabla \times \nabla \phi)_z = \partial_x (\partial_y \phi) - \partial_y (\partial_x \phi) = \phi_{yx} - \phi_{xy} = 0$. Geometrically, this implies that a gradient field is "irrotational"; it represents a "conservative" force field where the path integral between two points is independent of the route taken, meaning it cannot "circulate" around a point.}:
\begin{align*}
    \nabla \times \vec{E} &= \nabla \times \left( -\nabla \phi - \frac{1}{c}\frac{\partial \vec{A}}{\partial t} \right) \\
    &= -\nabla \times (\nabla \phi) - \frac{1}{c}\frac{\partial}{\partial t} (\nabla \times \vec{A}) \\
    &= 0 - \frac{1}{c}\frac{\partial \vec{B}}{\partial t} = -\frac{1}{c}\frac{\partial \vec{B}}{\partial t}
\end{align*}
Both equations are satisfied completely by the mathematical definition of the potentials.

\textbf{Step 3: Derive wave equations using the Lorentz gauge.}
Before performing the derivation, we must introduce the \textbf{Lorentz Gauge Condition}. As established in Step 1, we have the "gauge freedom" to add a gradient to $\vec{A}$ and a time-derivative to $\phi$ without changing the physical fields $\vec{E}$ and $\vec{B}$. The Lorentz gauge is a specific choice of this freedom, defined by the constraint:
\begin{equation*}
    \frac{1}{c} \frac{\partial \phi}{\partial t} + \nabla \cdot \vec{A} = 0
\end{equation*}
We choose this specific condition for two primary reasons:
\begin{enumerate}
    \item \textbf{Decoupling:} In the general Maxwell equations, the equations for $\phi$ and $\vec{A}$ are "coupled" (one depends on the other). Applying the Lorentz gauge eliminates the cross-terms, allowing each potential to be solved as an independent wave equation.
    \item \textbf{Lorentz Covariance:} Unlike other choices (such as the Coulomb gauge), the Lorentz gauge is "manifestly covariant." This means it treats space and time symmetrically, ensuring that the gauge condition looks the same to all observers moving at different speeds (inertial frames).
\end{enumerate}

Substituting the potentials into the inhomogeneous Maxwell equations (Gauss's Law and Ampere-Maxwell's Law):
\begin{equation*}
    \nabla \cdot \vec{E} = 4\pi\rho, \quad \nabla \times \vec{B} - \frac{1}{c}\frac{\partial \vec{E}}{\partial t} = \frac{4\pi}{c}\vec{j}
\end{equation*}
For Gauss's Law:
\begin{equation*}
    \nabla \cdot \left( -\nabla \phi - \frac{1}{c}\frac{\partial \vec{A}}{\partial t} \right) = -\nabla^2 \phi - \frac{1}{c}\frac{\partial}{\partial t} (\nabla \cdot \vec{A}) = 4\pi\rho
\end{equation*}
Applying the Lorentz gauge $\nabla \cdot \vec{A} = -\frac{1}{c} \partial_t \phi$:
\begin{equation*}
    -\nabla^2 \phi - \frac{1}{c}\frac{\partial}{\partial t} \left( -\frac{1}{c} \frac{\partial \phi}{\partial t} \right) = 4\pi\rho \implies \left( \frac{1}{c^2} \frac{\partial^2}{\partial t^2} - \nabla^2 \right) \phi = 4\pi\rho
\end{equation*}
Similarly, for Ampere-Maxwell's Law, using the identity $\nabla \times (\nabla \times \vec{A}) = \nabla(\nabla \cdot \vec{A}) - \nabla^2 \vec{A}$:
\begin{align*}
    \nabla(\nabla \cdot \vec{A}) - \nabla^2 \vec{A} - \frac{1}{c}\frac{\partial}{\partial t} \left( -\nabla \phi - \frac{1}{c}\frac{\partial \vec{A}}{\partial t} \right) &= \frac{4\pi}{c}\vec{j} \\
    \left( \frac{1}{c^2}\frac{\partial^2}{\partial t^2} - \nabla^2 \right) \vec{A} + \nabla \left( \nabla \cdot \vec{A} + \frac{1}{c}\frac{\partial \phi}{\partial t} \right) &= \frac{4\pi}{c}\vec{j}
\end{align*}
Applying the Lorentz gauge condition again, the term in the second parenthesis vanishes, leaving:
\begin{equation*}
    \left( \frac{1}{c^2} \frac{\partial^2}{\partial t^2} - \nabla^2 \right) \vec{A} = \frac{4\pi}{c}\vec{j}
\end{equation*}
Thus, both equations decouple into standard wave equations $\Box \phi = 4\pi\rho$ and $\Box \vec{A} = \frac{4\pi}{c}\vec{j}$.


\textbf{Step 4: Use the Green's function to compute potentials for a moving charge.}
For a point particle with charge $e$ at trajectory $\vec{R}(t)$, the charge and current densities are:
\begin{equation*}
    \rho(t', \vec{r}') = e \delta^3(\vec{r}' - \vec{R}(t')), \quad \vec{j}(t', \vec{r}') = e \vec{v}(t') \delta^3(\vec{r}' - \vec{R}(t'))
\end{equation*}
The solution using the Green's function is the convolution over all space and past times:
\begin{equation*}
    \phi(t, \vec{r}) = \int dt' d^3\vec{r}' G(t-t', \vec{r}-\vec{r}') (4\pi \rho(t', \vec{r}'))
\end{equation*}
Substituting the given Green's function $G = \frac{\theta(t-t')}{4\pi|\vec{r}-\vec{r}'|} \delta\left(t-t' - \frac{|\vec{r}-\vec{r}'|}{c}\right)$ into the integral for $\phi$:
\begin{equation*}
    \phi(t, \vec{r}) = \int_{-\infty}^{\infty} dt' \int_{\mathbb{R}^3} d^3\vec{r}' \frac{\delta\left(t-t' - \frac{|\vec{r}-\vec{r}'|}{c}\right)}{|\vec{r}-\vec{r}'|} e \delta^3(\vec{r}' - \vec{R}(t'))
\end{equation*}
To evaluate the inner 3D spatial integral, we use the \textbf{sifting property} of the Dirac delta function\sidenote{For any continuous function $f(\vec{r}')$, the integral over the entire space $\mathbb{R}^3$ with a delta function centered at $\vec{r}_0$ is simply the value of the function at that point: $\int_{\mathbb{R}^3} f(\vec{r}') \delta^3(\vec{r}' - \vec{r}_0) d^3\vec{r}' = f(\vec{r}_0)$. In this case, our "point of interest" is the particle's position at time $t'$, which is $\vec{r}_0 = \vec{R}(t')$.}. We treat the time $t'$ as a fixed parameter during the spatial integration over $\vec{r}'$:
\begin{align*}
    \phi(t, \vec{r}) &= e \int_{-\infty}^{\infty} dt' \left[ \int_{\mathbb{R}^3} d^3\vec{r}' \frac{\delta\left(t-t' - \frac{|\vec{r}-\vec{r}'|}{c}\right)}{|\vec{r}-\vec{r}'|} \delta^3(\vec{r}' - \vec{R}(t')) \right] \\
    &= e \int_{-\infty}^{\infty} dt' \left. \left( \frac{\delta\left(t-t' - \frac{|\vec{r}-\vec{r}'|}{c}\right)}{|\vec{r}-\vec{r}'|} \right) \right|_{\vec{r}' = \vec{R}(t')} \\
    &= \int_{-\infty}^{\infty} dt' \frac{e}{|\vec{r}-\vec{R}(t')|} \delta\left(t-t' - \frac{|\vec{r}-\vec{R}(t')|}{c}\right)
\end{align*}
This reduces the 4D spacetime integral to a 1D integral over the particle's history. The remaining delta function enforces that we only sum the contribution from the specific past time $t'$ when a signal traveling at speed $c$ from $\vec{R}(t')$ could have reached our current position $\vec{r}$ at the current time $t$.
Under the far-field non-relativistic approximation, $|\vec{R}(t')| \ll |\vec{r}|$ and $|\vec{v}| \ll c$. We can approximate the distance as simply $|\vec{r}-\vec{R}(t')| \approx |\vec{r}| = r$.
Thus, evaluating the delta function sets $t' = t - r/c$:
\begin{equation*}
    \phi(t, \vec{r}) \approx \frac{e}{r}
\end{equation*}
Similarly, for the vector potential:
\begin{align*}
    \vec{A}(t, \vec{r}) &= \frac{1}{c} \int dt' d^3\vec{r}' \frac{\delta\left(t-t' - \frac{|\vec{r}-\vec{r}'|}{c}\right)}{|\vec{r}-\vec{r}'|} e \vec{v}(t') \delta^3(\vec{r}' - \vec{R}(t')) \\
    &= \frac{e}{c} \int dt' \frac{\vec{v}(t')}{|\vec{r}-\vec{R}(t')|} \delta\left(t-t' - \frac{|\vec{r}-\vec{R}(t')|}{c}\right)
\end{align*}
Using the exact same far-field approximation, we pull out $1/r$ and evaluate at the retarded time $t_{ret} = t - r/c$:
\begin{equation*}
    \vec{A}(t, \vec{r}) \approx \frac{e}{cr} \vec{v}(t - r/c)
\end{equation*}
These represent the leading-order Liénard-Wiechert potentials for a slowly moving charge viewed from far away.
\end{solution}

\textbf{Remark: The Geometry of Gauge --- Equivalence Classes and Redundancy}

To a mathematician, a "Gauge" is best understood as an \textbf{equivalence class} within the configuration space of all possible potentials $\mathcal{A}$. The physical fields $\vec{E}$ and $\vec{B}$ do not uniquely determine a single potential $A = (\phi, \vec{A})$. Instead, we define an equivalence relation $\sim$ such that $A_1 \sim A_2$ if and only if they differ by the derivative of a scalar field $\chi$:
\begin{equation*}
    \vec{A}_2 = \vec{A}_1 + \nabla \chi, \quad \phi_2 = \phi_1 - \frac{1}{c} \frac{\partial \chi}{\partial t}
\end{equation*}
The collection of all such equivalent potentials forms a \textbf{gauge orbit} $[A]$. The "physical reality" is not the representative potential itself, but the orbit $[A]$. This structure is identical to the one found in \textbf{De Rham Cohomology}: since the field strength $F$ is a 2-form such that $F = dA$, then $A$ and $A+d\chi$ belong to the same cohomology class because $d(d\chi) = 0$.

\textbf{Why Fix a Gauge? Comparing Lorentz and Coulomb Choices}
Because of this redundancy, we have the freedom to impose an additional constraint—a process known as \textbf{Gauge Fixing}. This is mathematically analogous to choosing a specific coordinate system to simplify a geometry problem.

\begin{itemize}
    \item \textbf{Lorentz Gauge ($\nabla \cdot \vec{A} + \frac{1}{c} \partial_t \phi = 0$):} This is the choice made in this exercise. Its primary mathematical advantage is that it treats space and time on equal footing, making it \textbf{Lorentz invariant} (consistent with Einstein's Relativity). In this gauge, the coupled Maxwell equations decouple into independent wave equations:
    \begin{equation*}
        \Box \phi = 4\pi\rho, \quad \Box \vec{A} = \frac{4\pi}{c}\vec{j}
    \end{equation*}
    where $\Box$ is the d'Alembertian operator. This gauge is the "natural" choice for problems involving light, radiation, and causality.

    \item \textbf{Coulomb Gauge ($\nabla \cdot \vec{A} = 0$):} This choice, also called the "transverse gauge," focuses on making the scalar potential $\phi$ as simple as possible. Under this constraint, Gauss's law becomes a simple \textbf{Poisson equation}:
    \begin{equation*}
        \nabla^2 \phi = -4\pi\rho
    \end{equation*}
    This implies that $\phi$ reacts "instantly" to charge distributions. While this seems to contradict relativity, the vector potential $\vec{A}$ compensates for this in such a way that the physical fields $\vec{E}$ and $\vec{B}$ still respect the speed of light. This gauge is the preferred tool for atomic physics and quantum optics where static interactions are dominant.
\end{itemize}

Philosophically, gauge fixing tells us that while our \textit{mathematical descriptions} may contain redundant information, the \textit{laws of nature} are invariant under these changes. Choosing a gauge is simply the act of picking the most efficient language to describe an underlying geometric truth.

\subsection{Statistical Mechanics}

\begin{exercise}
\label{physics:2023:4}
Consider a one-dimensional system of free massless bosons with one polarization, and the dispersion relation $E _ { k } = \hbar v | k |$, where $v$ is the particle velocity, $k$ wavevector, $E_k$ energy. The particles are not interacting either among themselves or with external scattering potentials. If the system is in equilibrium at temperature $T$,

\begin{enumerate}[label=(\alph*)]
    \item Assume that the chemical potential $\mu$ is $0$ , calculate the heat capacity $C$ per unit length.
    \item Repeat the calculations for massive fermions, for which $E _ { k } = \hbar ^ { 2 } k ^ { 2 } / 2 m$ and the chemical potential $\mu$ is far above the bottom of the energy spectrum: $\mu \gg k_B T$. (Consider one spin direction for the fermions, and you may make reasonable approximations.)
\end{enumerate}

Hint: you might find the following formulae useful to evaluate some integrals needed
\begin{equation}
    \sum_{k=1}^\infty \frac{1}{k^2} = \frac{\pi^2}{6}, \quad \sum_{k=1}^\infty \frac{(-1)^{k+1}}{k^2} = \frac{\pi^2}{12}. \label{eq:2023:4:8}
\end{equation}
\end{exercise}

\begin{solution}
\textbf{Prerequisites / Core Concepts:}
\begin{itemize}
    \item \textbf{Internal Energy:} The total energy of a system, found by summing the energy of each state multiplied by its average occupation number: $U = \int E g(E) n(E) dE$.
    \item \textbf{Bose-Einstein Statistics:} The distribution function for identical bosons: $n_B(E) = \frac{1}{e^{(E-\mu)/(k_B T)} - 1}$.
    \item \textbf{Fermi-Dirac Statistics:} The distribution function for identical fermions: $n_F(E) = \frac{1}{e^{(E-\mu)/(k_B T)} + 1}$.
    \item \textbf{Density of States:} The number of available quantum states per unit energy $g(E)$, defined as $g(E) = \frac{L}{2\pi} \frac{dk}{dE}$ per dimension, multiplied by degeneracy factors.
\end{itemize}

\textbf{Intuition / Motivation:}
Think of a stadium with seats arranged by height (energy). For bosons, everyone can crowd into the lowest seats. At low temperatures, only a few get bumped up slightly. For fermions, only one person can sit in each seat, so they stack up high. At low temperatures, only the few people right at the top surface can jump higher. The heat capacity measures how much energy it takes to "bump" people up. Because the stadium is 1D (a narrow aisle), the number of available seats at a given height behaves differently for massless particles (constant seats) versus massive ones (fewer seats as you go up), which directly dictates how the heat capacity scales with temperature.

\textbf{Logical Step-by-Step Breakdown:}
\textbf{Step 1: Calculate the internal energy for 1D massless bosons.}
We are given a 1D system of length $L$. The wavevectors are quantized as $k = \frac{2\pi n}{L}$. The sum over states converts to an integral: $\sum_k \to \frac{L}{2\pi} \int_{-\infty}^{\infty} dk$. Since the energy $E_k = \hbar v |k|$ depends only on the magnitude $|k|$, we can integrate from $0$ to $\infty$ and multiply by 2 (to account for $+k$ and $-k$).
The average occupation number for bosons with $\mu = 0$ is $n_B(k) = \frac{1}{e^{\beta \hbar v k} - 1}$, where $\beta = \frac{1}{k_B T}$.
The total internal energy $U$ is:
\begin{equation*}
    U = 2 \int_0^\infty \frac{L}{2\pi} dk \, E_k n_B(k) = \frac{L}{\pi} \int_0^\infty dk \frac{\hbar v k}{e^{\beta \hbar v k} - 1}
\end{equation*}

\textbf{Step 2: Evaluate the integral and heat capacity for bosons.}
To solve the integral, we introduce the dimensionless variable $x = \beta \hbar v k$. Then $dk = \frac{k_B T}{\hbar v} dx$ and $k = \frac{k_B T}{\hbar v} x$.
Substituting this into the integral:
\begin{equation*}
    U = \frac{L \hbar v}{\pi} \int_0^\infty \left( \frac{k_B T}{\hbar v} dx \right) \left( \frac{k_B T}{\hbar v} x \right) \frac{1}{e^x - 1} = \frac{L (k_B T)^2}{\pi \hbar v} \int_0^\infty \frac{x}{e^x - 1} dx
\end{equation*}
We are given the standard integral result $\int_0^\infty \frac{x}{e^x - 1} dx = \frac{\pi^2}{6}$. Thus:
\begin{equation*}
    U = \frac{L (k_B T)^2}{\pi \hbar v} \left( \frac{\pi^2}{6} \right) = \frac{\pi L k_B^2 T^2}{6 \hbar v}
\end{equation*}
The heat capacity per unit length $C/L$ is the derivative of the internal energy with respect to temperature:
\begin{equation*}
    \frac{C}{L} = \frac{1}{L} \frac{\partial U}{\partial T} = \frac{\pi k_B^2}{6 \hbar v} (2T) = \frac{\pi k_B^2 T}{3 \hbar v}
\end{equation*}

\textbf{Step 3: Calculate the density of states for 1D massive fermions.}
For massive fermions, the dispersion relation is $E_k = \frac{\hbar^2 k^2}{2m}$. We only consider $k \ge 0$ with a factor of 2 for $\pm k$. The density of states (considering only one spin direction, so no spin degeneracy factor) is:
\begin{equation*}
    g(E) dE = 2 \frac{L}{2\pi} dk = \frac{L}{\pi} \frac{dk}{dE} dE
\end{equation*}
From $E = \frac{\hbar^2 k^2}{2m}$, we have $k = \frac{\sqrt{2mE}}{\hbar}$, so $\frac{dk}{dE} = \frac{\sqrt{2m}}{\hbar} \frac{1}{2\sqrt{E}} = \sqrt{\frac{m}{2\hbar^2 E}}$.
Thus, the density of states is:
\begin{equation*}
    g(E) = \frac{L}{\pi \hbar} \sqrt{\frac{m}{2E}}
\end{equation*}

\textbf{Step 4: Use the Sommerfeld expansion for the fermion heat capacity.}
For a highly degenerate Fermi gas where $\mu \gg k_B T$, the chemical potential $\mu$ is approximately the Fermi energy $E_F$. Only fermions within a narrow energy shell of width $\sim k_B T$ around the Fermi surface can be thermally excited.
The Sommerfeld expansion gives a general formula for the low-temperature heat capacity of a Fermi gas:
\begin{equation*}
    C = \frac{\pi^2}{3} k_B^2 T g(\mu)
\end{equation*}
Substituting our density of states evaluated at the chemical potential $\mu$:
\begin{equation*}
    C = \frac{\pi^2}{3} k_B^2 T \left( \frac{L}{\pi \hbar} \sqrt{\frac{m}{2\mu}} \right) = \frac{\pi L k_B^2 T}{3\hbar} \sqrt{\frac{m}{2\mu}}
\end{equation*}
Therefore, the heat capacity per unit length is:
\begin{equation*}
    \frac{C}{L} = \frac{\pi k_B^2 T}{3\hbar} \sqrt{\frac{m}{2\mu}}
\end{equation*}
In both cases (massless bosons and massive fermions), we see that the heat capacity per unit length in 1D is strictly proportional to the temperature $T$.
\end{solution}

\subsection{General Relativity}

\begin{exercise}
\label{physics:2023:5}
Consider the vacuum Einstein's equation in four dimensional spacetime with a cosmological constant
\begin{equation}
    R_{\mu\nu} - \frac{1}{2} g_{\mu\nu} R + g_{\mu\nu} \Lambda = 0. \label{eq:2023:5:9}
\end{equation}

\begin{enumerate}[label=(\alph*)]
    \item Prove that $R _ { \mu \nu } = k g _ { \mu \nu }$ and find out the value of $k$ .
    \item Now start with an ansatz of a metric in the following form
    \begin{equation}
        ds^2 = -f(r) dt^2 + \frac{1}{f(r)} dr^2 + r^2 (d\theta^2 + \sin^2\theta d\phi^2), \label{eq:2023:5:10}
    \end{equation}
    where $f(r)$ is a polynomial in $r$. Compute non zero components of the Ricci tensor $R_{\mu\nu}$ and scalar curvature $R$ of this metric.
    \item Assuming that the above ansatz is a solution of the vacuum Einstein equation with cosmological constant $\Lambda$, solve $f(r)$.
    \item Prove that $\partial _ { t }$ and $\partial _ { \phi }$ are Killing vector fields.
\end{enumerate}
\end{exercise}

\begin{solution}
\textbf{Prerequisites / Core Concepts:}
\begin{itemize}
    \item \textbf{The Riemann Curvature Tensor ($R^\rho_{\sigma\mu\nu}$):} The most fundamental measure of curvature. It describes how a vector changes when parallel-transported around an infinitesimal closed loop. In 4D, it has 256 components (though only 20 are independent due to symmetries).
    \item \textbf{Ricci Tensor ($R_{\mu\nu}$):} A "contraction" (a generalized trace) of the Riemann tensor: $R_{\mu\nu} = R^\alpha_{\mu\alpha\nu}$. While the Riemann tensor tracks all distortion, the Ricci tensor specifically measures the change in the \textbf{volume} of a small region of space as it moves along a geodesic.
    \item \textbf{Scalar Curvature ($R$):} The trace of the Ricci tensor with respect to the metric: $R = g^{\mu\nu} R_{\mu\nu}$. It is a single scalar value at each point in spacetime that represents the "average" curvature.
    \item \textbf{Einstein Field Equations (EFE):} The set of 10 coupled, non-linear partial differential equations that relate the geometry of spacetime (left side) to the energy-momentum content (right side). In vacuum with a cosmological constant $\Lambda$, they take the form $R_{\mu\nu} - \frac{1}{2} R g_{\mu\nu} + \Lambda g_{\mu\nu} = 0$.
    \item \textbf{The Ansatz Method:} A problem-solving technique where one "guesses" a mathematical form for the solution based on physical symmetries (e.g., assuming a star is perfectly spherical and static) to reduce complex PDEs into solvable ODEs.
\end{itemize}

\textbf{Remark: Geometry and the Ansatz --- How to Solve Spacetime}

For a mathematics student, the transition from Calculus to General Relativity can be understood as moving from "functions on flat space" to "the geometry of the space itself."

\textbf{1. Ricci Tensor ($R_{\mu\nu}$) vs. Scalar Curvature ($R$)}
Think of the difference between a matrix and its trace. 
The \textbf{Ricci Tensor} is like a $4 \times 4$ matrix at every point in space. It captures how curvature affects different directions. Specifically, if you have a small ball of test particles, the Ricci tensor tells you how the \textit{volume} of that ball grows or shrinks as it falls through gravity.
The \textbf{Scalar Curvature} is a single number. It is the sum of the diagonal elements of the Ricci tensor (adjusted by the metric). It tells you how the volume of a sphere in curved space differs from the volume of a sphere in flat Euclidean space. 
\textit{Analogy:} If the Riemann tensor is the full high-definition video of the curvature, the Ricci tensor is a high-quality photograph, and the Scalar curvature is a single brightness value.

\textbf{2. What is an "Ansatz"?}
The word \textbf{Ansatz} (German for "approach" or "placement") is used when we have a set of equations that are too difficult to solve from scratch. Einstein's equations are 10 non-linear PDEs; solving them for a general case is impossible.
An \textbf{Ansatz} is a "proactive guess." By observing that a star is spherical and doesn't change over time, we \textit{assume} the metric must look like:
\begin{equation*}
    ds^2 = -f(r) dt^2 + \frac{1}{f(r)} dr^2 + r^2 d\Omega^2
\end{equation*}
This guess (the Ansatz) contains an unknown function $f(r)$. We plug this guess into the 10 scary PDEs. Because the guess matches the symmetry of the problem, 8 of the equations might become $0=0$, and the remaining 2 become simple 1D calculus problems (ODEs) that allow us to solve for $f(r)$.

\textbf{3. Vacuum Einstein Equations}
When we say "vacuum," we mean there is no matter ($\text{Energy-Momentum Tensor } T_{\mu\nu} = 0$). However, the \textbf{Cosmological Constant} $\Lambda$ acts like a "dark energy" that fills the vacuum. Part (a) of this exercise proves a beautiful geometric result: in such a vacuum, the Ricci tensor is \textit{proportional} to the metric itself ($R_{\mu\nu} = \Lambda g_{\mu\nu}$). This is the definition of an \textbf{Einstein Manifold} in differential geometry.

\textbf{Logical Step-by-Step Breakdown:}
\textbf{Step 1: Answer Part (a) --- Proving $R_{\mu\nu} = kg_{\mu\nu}$ and finding $k$.}
The vacuum Einstein field equations with a cosmological constant are:
\begin{equation}
    R_{\mu\nu} - \frac{1}{2} g_{\mu\nu} R + g_{\mu\nu} \Lambda = 0 \label{eq:efe_vac}
\end{equation}
To find the relation between $R_{\mu\nu}$ and $g_{\mu\nu}$, we take the trace of the entire equation by contracting it with the inverse metric $g^{\mu\nu}$:
\begin{align*}
    g^{\mu\nu} \left( R_{\mu\nu} - \frac{1}{2} g_{\mu\nu} R + g_{\mu\nu} \Lambda \right) &= 0 \\
    g^{\mu\nu} R_{\mu\nu} - \frac{1}{2} (g^{\mu\nu} g_{\mu\nu}) R + (g^{\mu\nu} g_{\mu\nu}) \Lambda &= 0
\end{align*}
In 4-dimensional spacetime, the trace of the Ricci tensor is $g^{\mu\nu} R_{\mu\nu} = R$ (the Ricci scalar) and the trace of the metric is $g^{\mu\nu} g_{\mu\nu} = \delta^\mu_\mu = 4$. Substituting these:
\begin{align*}
    R - \frac{1}{2}(4) R + 4\Lambda &= 0 \\
    R - 2R + 4\Lambda &= 0 \implies R = 4\Lambda
\end{align*}
Now, substitute the scalar curvature $R = 4\Lambda$ back into the original vacuum field equation \eqref{eq:efe_vac}:
\begin{align*}
    R_{\mu\nu} - \frac{1}{2} g_{\mu\nu} (4\Lambda) + g_{\mu\nu} \Lambda &= 0 \\
    R_{\mu\nu} - 2\Lambda g_{\mu\nu} + \Lambda g_{\mu\nu} &= 0 \\
    R_{\mu\nu} - \Lambda g_{\mu\nu} &= 0 \implies R_{\mu\nu} = \Lambda g_{\mu\nu}
\end{align*}
This confirms that $R_{\mu\nu} = k g_{\mu\nu}$ with the constant proportionality factor \textbf{$k = \Lambda$}.

\textbf{Step 2: Answer Part (b) --- Compute non-zero Ricci components and Scalar Curvature.}
To compute the curvature of the spacetime, we follow the standard tensor hierarchy starting from the metric $g_{\mu\nu}$:
\begin{enumerate}
    \item \textbf{Christoffel Symbols ($\Gamma$):} These are computed from the metric derivatives:
    \begin{equation*}
        \Gamma^\sigma_{\mu\nu} = \frac{1}{2} g^{\sigma\rho} (\partial_\mu g_{\nu\rho} + \partial_\nu g_{\mu\rho} - \partial_\rho g_{\mu\nu})
    \end{equation*}
    \item \textbf{Riemann Curvature Tensor ($R$):} This measures the non-commutativity of covariant derivatives:
    \begin{equation*}
        R^\rho_{\sigma\mu\nu} = \partial_\mu \Gamma^\rho_{\nu\sigma} - \partial_\nu \Gamma^\rho_{\mu\sigma} + \Gamma^\rho_{\mu\lambda} \Gamma^\lambda_{\nu\sigma} - \Gamma^\rho_{\nu\lambda} \Gamma^\lambda_{\mu\sigma}
    \end{equation*}
    \item \textbf{Ricci Tensor ($R_{\mu\nu}$):} This is the specific contraction of the Riemann tensor over the first and third indices:
    \begin{equation*}
        R_{\sigma\nu} = R^\mu_{\sigma\mu\nu} = \partial_\mu \Gamma^\mu_{\nu\sigma} - \partial_\nu \Gamma^\mu_{\mu\sigma} + \Gamma^\mu_{\mu\lambda} \Gamma^\lambda_{\nu\sigma} - \Gamma^\mu_{\nu\lambda} \Gamma^\lambda_{\mu\sigma}
    \end{equation*}
\end{enumerate}

\textbf{Remark: Mnemonic for Curvature Tensors and the Logic of Contraction}

If you are a math student facing these formulas in an exam, do not try to memorize every index individually. Instead, memorize the \textbf{structural pattern}:
\begin{enumerate}
    \item \textbf{The "Derivative-Quadratic" Pattern:} Both the Riemann and Ricci tensors follow the same rhythmic structure:
    \begin{equation*}
        \text{Tensor} = (\partial \Gamma - \partial \Gamma) + (\Gamma\Gamma - \Gamma\Gamma)
    \end{equation*}
    Think of it as "Difference of derivatives" plus "Difference of products." 
    \item \textbf{The Index Flow:} In the Riemann tensor $R^\rho_{\sigma\mu\nu}$, the last two indices $(\mu, \nu)$ are the ones being derived ($\partial_\mu$ and $\partial_\nu$). The negative signs always fall on the term where $\mu$ and $\nu$ are swapped.
    \item \textbf{The Ricci "Shortcut":} You asked if the Ricci tensor is just "replacing the role by $\mu$." \textbf{Exactly.} The Ricci tensor $R_{\sigma\nu}$ is the "trace" of the Riemann tensor. You simply set the top index $\rho$ and the first "lower-right" index $\mu$ to the same letter (say, $\lambda$) and sum over them:
    \begin{equation*}
        R_{\sigma\nu} = \sum_{\lambda} R^\lambda_{\sigma\lambda\nu}
    \end{equation*}
    In an exam, if you can write down the 4-index Riemann tensor, you can "collapse" it into the 2-index Ricci tensor just by matching those two indices.
\end{enumerate}
\textit{Analogy:} If the Riemann tensor is a $4 \times 4 \times 4 \times 4$ hypercube of information, the Ricci tensor is the result of "squashing" that hypercube into a $4 \times 4$ flat matrix. You lose some details (the shape distortion), but you keep the most important part (the volume change).

For the given metric ansatz $g_{tt} = -f$, $g_{rr} = f^{-1}$, $g_{\theta\theta} = r^2$, and $g_{\phi\phi} = r^2 \sin^2\theta$, the non-zero Christoffel symbols are derived using Step 1 above\sidenote{The full set of non-zero symbols for this metric are: $\Gamma^t_{tr} = f'/2f$, $\Gamma^r_{tt} = ff'/2$, $\Gamma^r_{rr} = -f'/2f$, $\Gamma^r_{\theta\theta} = -rf$, $\Gamma^r_{\phi\phi} = -rf\sin^2\theta$, $\Gamma^\theta_{r\theta} = 1/r$, $\Gamma^\theta_{\phi\phi} = -\sin\theta\cos\theta$, $\Gamma^\phi_{r\phi} = 1/r$, and $\Gamma^\phi_{\theta\phi} = \cot\theta$.}. 

Substituting these symbols into the general Ricci formula for each component:
\begin{itemize}
    \item $R_{tt} = \frac{f f''}{2} + \frac{f f'}{r}$
    \item $R_{rr} = -\frac{f''}{2f} - \frac{f'}{rf}$
    \item $R_{\theta\theta} = 1 - f - rf' = 1 - (rf)'$
    \item $R_{\phi\phi} = (1 - f - rf') \sin^2 \theta = R_{\theta\theta} \sin^2 \theta$
\end{itemize}
Now, we compute the \textbf{Scalar Curvature $R$} by contracting the Ricci tensor with the inverse metric $g^{\mu\nu}$:
\begin{align*}
    R &= g^{tt} R_{tt} + g^{rr} R_{rr} + g^{\theta\theta} R_{\theta\theta} + g^{\phi\phi} R_{\phi\phi} \\
      &= \left(-\frac{1}{f}\right) \left(\frac{ff''}{2} + \frac{ff'}{r}\right) + (f) \left(-\frac{f''}{2f} - \frac{f'}{rf}\right) + \frac{1}{r^2}(1-f-rf') + \frac{1}{r^2\sin^2\theta}(1-f-rf')\sin^2\theta \\
      &= -f'' - \frac{4f'}{r} + \frac{2(1-f)}{r^2}
\end{align*}
This completes the required computations for the Ricci components and scalar curvature.

\textbf{Step 3: Answer Part (c) --- Solve for $f(r)$.}
Assuming the ansatz is a solution to the vacuum Einstein equation with cosmological constant $\Lambda$, we apply the result from Part (a): $R_{\mu\nu} = \Lambda g_{\mu\nu}$.
For the $\theta\theta$ component, we equate the expression from Part (b) to the metric component:
\begin{equation*}
    R_{\theta\theta} = \Lambda g_{\theta\theta} \implies 1 - \frac{d}{dr} (rf) = \Lambda r^2
\end{equation*}
Rearranging to isolate the derivative:
\begin{equation*}
    \frac{d}{dr} (rf) = 1 - \Lambda r^2
\end{equation*}
Integrating both sides with respect to $r$:
\begin{align*}
    \int d(rf) &= \int (1 - \Lambda r^2) dr \\
    rf(r) &= r - \frac{\Lambda r^3}{3} + C
\end{align*}
Dividing by $r$, we obtain the general solution:
\begin{equation*}
    f(r) = 1 - \frac{\Lambda r^2}{3} + \frac{C}{r}
\end{equation*}
By comparing this to the Newtonian limit (where $f(r) \approx 1 - 2M/r$ for small $\Lambda$), we identify the integration constant $C = -2M$, where $M$ is the mass of the central object. Thus:
\begin{equation*}
    f(r) = 1 - \frac{2M}{r} - \frac{\Lambda r^2}{3}
\end{equation*}

\textbf{Step 4: Answer Part (d) --- Prove $\partial_t$ and $\partial_\phi$ are Killing vector fields.}
To rigorously prove that a vector field $\xi$ is a Killing vector field, we must show it satisfies Killing's equation $\mathcal{L}_\xi g_{\mu\nu} = 0$. Using the coordinate expression for the Lie derivative of a $(0,2)$ tensor:
\begin{equation*}
    \mathcal{L}_\xi g_{\mu\nu} = \xi^\alpha \partial_\alpha g_{\mu\nu} + g_{\alpha\nu} \partial_\mu \xi^\alpha + g_{\mu\alpha} \partial_\nu \xi^\alpha = 0
\end{equation*}

\begin{itemize}
    \item \textbf{Proof for $\xi = \partial_t$:} In this basis, the components of the vector field are $\xi^\alpha = \delta^\alpha_t$ (meaning $\xi^t = 1$ and all other components are zero).
    Substituting this into the Lie derivative formula:
    \begin{align*}
        \mathcal{L}_{\partial_t} g_{\mu\nu} &= (1) \partial_t g_{\mu\nu} + g_{\alpha\nu} \partial_\mu (\delta^\alpha_t) + g_{\mu\alpha} \partial_\nu (\delta^\alpha_t) \\
        &= \partial_t g_{\mu\nu} + 0 + 0
    \end{align*}
    since the derivatives of the constant Kronecker delta $\delta^\alpha_t$ vanish. From the metric $ds^2$, we see that all $g_{\mu\nu}$ components depend only on $r$ and $\theta$. Thus, $\partial_t g_{\mu\nu} \equiv 0$ for all $\mu, \nu$, which proves $\mathcal{L}_{\partial_t} g_{\mu\nu} = 0$.

    \item \textbf{Proof for $\xi = \partial_\phi$:} Here, the components are $\xi^\alpha = \delta^\alpha_\phi$.
    The Lie derivative becomes:
    \begin{align*}
        \mathcal{L}_{\partial_\phi} g_{\mu\nu} &= (1) \partial_\phi g_{\mu\nu} + g_{\alpha\nu} \partial_\mu (\delta^\alpha_\phi) + g_{\mu\alpha} \partial_\nu (\delta^\alpha_\phi) \\
        &= \partial_\phi g_{\mu\nu} + 0 + 0
    \end{align*}
    By inspection of the metric, none of the components $g_{tt}(r)$, $g_{rr}(r)$, $g_{\theta\theta}(r)$, or $g_{\phi\phi}(r, \theta)$ contain the variable $\phi$. Therefore, $\partial_\phi g_{\mu\nu} \equiv 0$, confirming that $\mathcal{L}_{\partial_\phi} g_{\mu\nu} = 0$.
\end{itemize}
These rigorous calculations demonstrate that $\partial_t$ and $\partial_\phi$ generate isometries of the Schwarzschild-de Sitter spacetime, corresponding to its stationarity and axisymmetry.
\end{solution}

\subsection{Quantum Field Theory}

\begin{exercise}
\label{physics:2023:6}
Consider the $\phi ^ { 3 }$ model\sidenote{The name $\phi^3$ comes from the interaction term in the Lagrangian. In QFT, the power of the field dictates the 'topology' of the interactions: a $\phi^3$ term represents a vertex where three lines meet, signifying one particle splitting into two, or two merging into one.} with a real scalar field $\phi ( x )$ in $3 + 1$ dimensional Minkowski spacetime with metric $( - , + , + , + )$. Its Lagrangian density is
\begin{equation}
    \mathcal {L} = - \frac {1}{2} \partial_ {\mu} \phi \partial^ {\mu} \phi - \frac {1}{2} m ^ {2} \phi^ {2} - \frac {1}{6} g \phi^ {3}, \label{eq:2023:6:11}
\end{equation}
where $g$ is a coupling with dimensions of mass.

\textbf{Remark: The Role of the $\phi^3$ Toy Model}
To a mathematician, Quantum Field Theory (QFT) is the study of functional integrals over a space of fields. The $\phi^3$ model is what physicists call a \textbf{"Toy Model"}.
\begin{itemize}
    \item \textbf{Why use it?}: It is the simplest possible field theory that includes interactions. The linear and quadratic terms (the "Free Theory") describe particles that never see each other. The $\phi^3$ term introduces the first "drama"---it allows particles to scatter, decay, and produce "virtual" particle loops.
    \item \textbf{Mathematical Feature}: It is the perfect laboratory for learning \textbf{Renormalization}. Because the interactions are so simple, the infinite integrals (loops) that arise are easy to classify and "tame" using the dimensional regularization techniques shown in this exercise.
    \item \textbf{The Catch}: Physically, the $\phi^3$ model is unstable. Because the potential $V(\phi) = \frac{1}{2}m^2\phi^2 + \frac{1}{6}g\phi^3$ goes to $-\infty$ as $\phi \to -\infty$, the system has no true "bottom" or ground state. In real-world physics, we usually use the $\phi^4$ model (which is stable) or Gauge theories (which we saw in Exercise 3).
\end{itemize}

\begin{enumerate}[label=(\alph*)]
    \item Write down the propagator and the interaction vertex for this model in momentum space.
    \item Compute the one-loop self-energy graph using dimensional regularization.
    \item Introducing $m ^ { 2 } = m _ { R } ^ { 2 } + \delta m ^ { 2 }$. What is the value of $\delta m ^ { 2 }$ if we want to write the one-loop self-energy graph as a finite function of $m _ { R }$?
\end{enumerate}

You may find the following formula useful:
\begin{align}
    (A B) ^ {- 1} &= \int_ {0} ^ {1} d x [ x A + (1 - x) B ] ^ {- 2}, \label{eq:2023:6:12} \\
    \int d ^ {d} k \frac {1}{(- k ^ {2} - 2 p \cdot k - M ^ {2} + i \epsilon) ^ {s}} &= (- 1) ^ {s} i \pi^ {d / 2} \frac {\Gamma (s - d / 2)}{\Gamma (s)} (- p ^ {2} + M ^ {2} - i \epsilon) ^ {d / 2 - s}, \label{eq:2023:6:13}
\end{align}
where $\Gamma ( z )$ is the Gamma function which has a simple pole at the origin:
\begin{equation}
    \Gamma(z) = \frac{1}{z} - \gamma + \mathcal{O}(z), \label{eq:2023:6:14}
\end{equation}
where $\gamma$ is the Euler constant.
\end{exercise}

\textbf{Remark: The Landscape of Quantum Field Theory --- From Particles to Fields}

To a mathematician, Quantum Field Theory (QFT) is the result of applying the principles of Quantum Mechanics to systems with an infinite number of degrees of freedom.

\begin{enumerate}
    \item \textbf{The Necessity of Fields:} In ordinary Quantum Mechanics (Exercise 2), we track a fixed number of particles. However, Einstein's $E=mc^2$ tells us that energy can be converted into mass. To describe processes where particles are created or destroyed (like in a particle collider), we cannot use a single-particle wavefunction. Instead, we use a \textbf{Field} $\phi(x)$, which is defined at every point in spacetime.
    \item \textbf{The "Harmonic Oscillator" Analogy:} Think of spacetime as a fine mattress made of an infinite number of tiny springs. Every point in space has its own "spring" (a harmonic oscillator). 
    \begin{itemize}
        \item A \textbf{"Free Field"} ($\phi^2$ terms) means the springs are vibrating but not touching each other. These vibrations are what we perceive as "particles."
        \item An \textbf{"Interaction"} ($\phi^3$ terms) means the springs are linked. When one spring vibrates, it can transfer energy to its neighbors, causing them to vibrate too. This is how particles "scatter" or "decay."
    \end{itemize}
    \item \textbf{The Path Integral Formulation:} We compute "amplitudes" (probabilities) by summing over \textit{all possible ways} a field can evolve. This is a functional integral $\int \mathcal{D}\phi \, e^{iS[\phi]}$, where $S$ is the Action. Since we cannot solve this integral exactly for interacting theories, we use \textbf{Perturbation Theory} (Feynman diagrams) to approximate it, treating the interaction $g$ as a small parameter.
\end{enumerate}

\begin{solution}
\textbf{Prerequisites / Core Concepts:}
\begin{itemize}
    \item \textbf{Lagrangian Density ($\mathcal{L}$):} A local function $\mathcal{L}(\phi, \partial_\mu \phi)$ such that the Action $S = \int d^4x \, \mathcal{L}$ is minimized. It defines the dynamics of the field at every point in spacetime.
    \item \textbf{The Scalar Field ($\phi$):} A map $\phi: \mathbb{M}^4 \to \mathbb{R}$ from Minkowski spacetime to the real numbers. It represents particles with "spin 0" (like the Higgs boson).
    \item \textbf{The Propagator ($\Delta$):} Mathematically, this is the \textbf{Green's function} of the differential operator $(\Box + m^2)$. It represents the probability amplitude for a particle to travel from point $A$ to point $B$. In momentum space, it is the algebraic inverse of the kinetic operator: $\frac{-i}{p^2+m^2}$.
    \item \textbf{Interaction Vertex ($g$):} The point where the "drama" happens. A $\phi^3$ vertex is a point where three particle paths meet. The constant $g$ (coupling) measures the "strength" of this interaction.
    \item \textbf{Loop Diagrams and Self-Energy ($\Pi$):} A particle can emit a virtual particle and then quickly reabsorb it, forming a "loop" in its path. This process changes the particle's effective inertia, which we call \textbf{Self-Energy}.
    \item \textbf{Dimensional Regularization ($d = 4-\epsilon$):} Many QFT integrals are divergent (infinite). We "regularize" them by assuming we live in $4-\epsilon$ dimensions. This turns the infinity into a well-defined mathematical pole $1/\epsilon$.
    \item \textbf{Renormalization and Counterterms:} The "Bare" mass $m$ in our Lagrangian is not what we measure in a lab. We define a "Renormalized" mass $m_R$ and a "Counterterm" $\delta m^2$ such that $m^2 = m_R^2 + \delta m^2$. We then use $\delta m^2$ to "cancel out" the $1/\epsilon$ infinity, leaving a finite result for $m_R$.
\end{itemize}

\textbf{Intuition / Motivation: The "Dressed" Particle}
In the quantum world, a particle is never "naked." It is always surrounded by a "cloud" of virtual particles that it is constantly emitting and reabsorbing. 
Imagine a person (the bare particle) walking through a crowded room. As they move, they interact with people (virtual particles), which slows them down and makes them "feel" heavier. What we measure in the lab is the \textbf{"Dressed Particle"} (the person + the crowd). 
When we calculate this "cloud's" effect, we get an infinite result because there are infinitely many ways a particle can interact with high-energy virtual states. \textbf{Renormalization} is the mathematical procedure of redefining our parameters so that the infinite "cloud" effect is absorbed into the definition of the mass, leaving us with a finite, measurable "dressed" mass.

\textbf{Logical Step-by-Step Breakdown:}

\textbf{Step 1: Extract the Feynman rules from the Lagrangian.}
The Lagrangian density $\mathcal{L}$ for the $\phi^3$ model is:
\begin{equation*}
    \mathcal{L} = -\frac{1}{2} \partial_\mu \phi \partial^\mu \phi - \frac{1}{2} m^2 \phi^2 - \frac{1}{6} g \phi^3
\end{equation*}
Using the principle of least action, we derive the Feynman rules in momentum space:
\begin{itemize}
    \item \textbf{Propagator:} We isolate the "Free" part of the Lagrangian. Using integration by parts, the kinetic term $-\frac{1}{2}\partial_\mu \phi \partial^\mu \phi$ becomes $\frac{1}{2}\phi \Box \phi$. The free equation of motion is $(\Box + m^2)\phi = 0$. In momentum space, $\partial_\mu \to ip_\mu$, so $\Box \to -p^2$. The operator becomes $(-p^2 + m^2)$. The propagator $\Delta(p)$ is the $i$-multiplied inverse of this operator:
    \begin{equation*}
        \Delta(p) = \frac{-i}{p^2 + m^2}
    \end{equation*}
    \item \textbf{Interaction Vertex:} The interaction term is $-\frac{1}{3!} g \phi^3$. The vertex factor is $i$ times the coefficient of the interaction term, multiplied by the symmetry factor of the fields ($3!$). This yields a vertex factor of \textbf{$-ig$}.
\end{itemize}

\textbf{Step 2: Construct the one-loop self-energy integral.}
The self-energy diagram involves a particle of momentum $p$ splitting into two virtual particles with momenta $k$ and $p-k$, which then recombine.
\begin{itemize}
    \item \textbf{Symmetry Factor:} There are two internal lines representing the same boson. Swapping them does not create a new state, so we must divide the amplitude by $S = 2$.
    \item \textbf{The Setup:} Using the rules from Step 1, the amplitude $-i\Pi(p^2)$ is:
\end{itemize}
\begin{align*}
    -i\Pi(p^2) &= \frac{1}{2} (-ig)^2 \int \frac{d^d k}{(2\pi)^d} \Delta(k) \Delta(p-k) \\
    &= \frac{-g^2}{2} \int \frac{d^d k}{(2\pi)^d} \left( \frac{-i}{k^2 + m^2} \right) \left( \frac{-i}{(p-k)^2 + m^2} \right) \\
    &= \frac{i^2 g^2}{2} \int \frac{d^d k}{(2\pi)^d} \frac{1}{(k^2+m^2)[(p-k)^2+m^2]}
\end{align*}
Simplifying the $i^2 = -1$ and dividing by $-i$, we get:
\begin{equation*}
    \Pi(p^2) = \frac{i g^2}{2} \int \frac{d^d k}{(2\pi)^d} \frac{1}{(k^2+m^2)[(p-k)^2+m^2]}
\end{equation*}

\textbf{Step 3: Combine denominators using Feynman parameters.}
We use the mathematical trick $\frac{1}{AB} = \int_0^1 \frac{dx}{[xA + (1-x)B]^2}$. Let $A = (p-k)^2 + m^2$ and $B = k^2 + m^2$.
The combined denominator $D$ is:
\begin{align*}
    D &= x[(p-k)^2 + m^2] + (1-x)[k^2 + m^2] \\
      &= x[p^2 - 2p \cdot k + k^2 + m^2] + k^2 - xk^2 + m^2 - xm^2 \\
      &= k^2 - 2x(p \cdot k) + xp^2 + m^2
\end{align*}
To evaluate the integral over $k$, we "complete the square" by defining a shifted momentum $\ell = k - xp$.
Then $k = \ell + xp$, and:
\begin{align*}
    k^2 - 2xp \cdot k &= (\ell + xp)^2 - 2xp \cdot (\ell + xp) \\
    &= \ell^2 + 2x\ell \cdot p + x^2p^2 - 2x\ell \cdot p - 2x^2p^2 \\
    &= \ell^2 - x^2p^2
\end{align*}
Substituting this back into $D$:
\begin{equation*}
    D = \ell^2 - x^2p^2 + xp^2 + m^2 = \ell^2 + x(1-x)p^2 + m^2
\end{equation*}
Let $\Delta = x(1-x)p^2 + m^2$. The integral becomes:
\begin{equation*}
    \Pi(p^2) = \frac{i g^2}{2} \int_0^1 dx \int \frac{d^d \ell}{(2\pi)^d} \frac{1}{[\ell^2 + \Delta]^2}
\end{equation*}

\textbf{Step 4: Evaluate the momentum integral via Dimensional Regularization.}
We use the master formula $\int \frac{d^d \ell}{(2\pi)^d} \frac{1}{(\ell^2 + \Delta)^s} = \frac{i}{(4\pi)^{d/2}} \frac{\Gamma(s-d/2)}{\Gamma(s)} \Delta^{d/2-s}$.
Setting $s=2$ and $d=4-\epsilon$:
\begin{equation*}
    \int \frac{d^d \ell}{(2\pi)^d} \frac{1}{(\ell^2 + \Delta)^2} = \frac{i}{(4\pi)^{2-\epsilon/2}} \frac{\Gamma(\epsilon/2)}{\Gamma(2)} \Delta^{-\epsilon/2}
\end{equation*}
Now we perform the epsilon expansion ($\epsilon \to 0$):
\begin{itemize}
    \item $\frac{1}{(4\pi)^{2-\epsilon/2}} = \frac{1}{16\pi^2} (4\pi)^{\epsilon/2} \approx \frac{1}{16\pi^2} (1 + \frac{\epsilon}{2} \ln(4\pi))$.
    \item $\Gamma(\epsilon/2) = \frac{2}{\epsilon} - \gamma + \mathcal{O}(\epsilon)$.
    \item $\Delta^{-\epsilon/2} = e^{-\frac{\epsilon}{2}\ln\Delta} \approx 1 - \frac{\epsilon}{2}\ln\Delta$.
\end{itemize}
Multiplying these together and keeping only terms of order $1/\epsilon$ and $\mathcal{O}(1)$:
\begin{align*}
    \text{Integral} &\approx \frac{i}{16\pi^2} \left( \frac{2}{\epsilon} - \gamma + \ln(4\pi) - \ln\Delta \right) \\
    &= \frac{i}{16\pi^2} \left( \frac{2}{\epsilon} - \gamma + \ln \left( \frac{4\pi}{x(1-x)p^2 + m^2} \right) \right)
\end{align*}
Substituting this into our expression for $\Pi(p^2)$:
\begin{equation*}
    \Pi(p^2) \approx \frac{i g^2}{2} \int_0^1 dx \left[ \frac{i}{16\pi^2} \left( \frac{2}{\epsilon} - \dots \right) \right] = -\frac{g^2}{32\pi^2} \int_0^1 dx \left( \frac{2}{\epsilon} + \text{finite terms} \right)
\end{equation*}

\textbf{Step 5: Determine the mass counterterm $\delta m^2$.}
We define the "Renormalized" mass $m_R$ and "Counterterm" $\delta m^2$ such that the total mass squared is $m^2 = m_R^2 + \delta m^2$.
The physical self-energy must be finite. This requires that $\delta m^2$ cancels the divergent $1/\epsilon$ part of $\Pi(p^2)$.
From Step 4, the divergent part is:
\begin{equation*}
    \Pi_{div} = -\frac{g^2}{32\pi^2} \left( \frac{2}{\epsilon} \right) = -\frac{g^2}{16\pi^2 \epsilon}
\end{equation*}
To achieve cancellation, we set the counterterm to:
\begin{equation*}
    \delta m^2 = - \Pi_{div} = \frac{g^2}{16\pi^2 \epsilon}
\end{equation*}
This choice ensures that the final mass $m_R$ measured in experiments is a finite quantity, effectively absorbing the infinite effect of the virtual particle cloud into the unobservable bare parameter.
\end{solution}
